% !TeX program = xelatex
\documentclass[11pt,a4paper,oneside]{report}
\usepackage{fontspec}
\setmainfont{Times New Roman}
\setsansfont{Arial}
\setmonofont{Consolas}[Scale=.82]
\usepackage{polyglossia}
\setmainlanguage{vietnamese}
\usepackage[top=24mm,bottom=24mm,left=28mm,right=23mm,headheight=16pt]{geometry}
\usepackage{setspace}
\setstretch{1.16}
\usepackage{graphicx,array,booktabs,longtable,tabularx,amsmath,xcolor,enumitem}
\usepackage{titlesec,fancyhdr,caption,float,placeins,listings,needspace,microtype}
\usepackage{tikz}
\usetikzlibrary{arrows.meta,positioning,fit,calc}
\usepackage{hyperref,xurl}
\definecolor{vtx}{HTML}{EE0033}
\definecolor{navy}{HTML}{203340}
\definecolor{teal}{HTML}{006C73}
\definecolor{ViettelRed}{HTML}{EE0033}
\definecolor{pale}{HTML}{F1F6F7}
\definecolor{graytext}{HTML}{5B6670}
\hypersetup{colorlinks=true,linkcolor=navy,urlcolor=teal,citecolor=teal,
 pdftitle={Báo cáo tổng kết dự án - Mô phỏng hải chiến trên nền địa hình thực tế},
 pdfauthor={Lê Việt Thành; Nguyễn Hữu Hiệu},bookmarksnumbered=true,pdfdisplaydoctitle=true}
\urlstyle{tt}
\setlength{\parindent}{1.2em}
\setlength{\parskip}{4pt}
\setlength{\emergencystretch}{3em}
\setlength{\tabcolsep}{5pt}
\renewcommand{\arraystretch}{1.15}
\setlist[itemize]{leftmargin=1.5em,itemsep=2pt,topsep=3pt}
\setlist[enumerate]{leftmargin=1.8em,itemsep=3pt,topsep=3pt}
\captionsetup{font=small,labelfont=bf,labelsep=period}
\titleformat{\chapter}[display]{\sffamily\bfseries\color{navy}}{\large\color{vtx} CHƯƠNG \thechapter}{5pt}{\LARGE}
\titlespacing*{\chapter}{0pt}{-10pt}{17pt}
\titleformat{\section}{\large\sffamily\bfseries\color{navy}}{\thesection}{.65em}{}
\titleformat{\subsection}{\normalsize\sffamily\bfseries\color{teal}}{\thesubsection}{.65em}{}
\pagestyle{fancy}\fancyhf{}
\fancyhead[L]{\small\sffamily VIETTEL AEROSPACE}
\fancyhead[R]{\small\sffamily VDT 2026 / GIAI ĐOẠN 2}
\fancyfoot[L]{\footnotesize\color{graytext} Mô phỏng hải chiến trên địa hình thực tế}
\fancyfoot[R]{\thepage}
\renewcommand{\headrulewidth}{.4pt}
\fancypagestyle{plain}{\fancyhf{}\fancyfoot[R]{\thepage}\renewcommand{\headrulewidth}{0pt}}
\setcounter{tocdepth}{1}\setcounter{secnumdepth}{2}\raggedbottom
\newcommand{\note}[1]{\par\smallskip\noindent\fcolorbox{teal}{pale}{\parbox{\dimexpr\linewidth-2\fboxsep-2\fboxrule}{\small #1}}\par\smallskip}
\newcommand{\src}[1]{\par\noindent{\footnotesize\color{graytext}\textbf{Đối chiếu:} #1}\par}
\newcommand{\fronttitle}[1]{\begin{center}\Large\bfseries #1\end{center}\vspace{7mm}}
\lstset{basicstyle=\ttfamily\footnotesize,breaklines=true,frame=single,rulecolor=\color{teal!35},backgroundcolor=\color{pale},columns=fullflexible,keepspaces=true,showstringspaces=false,aboveskip=9pt,belowskip=9pt}
\tikzset{box/.style={draw=teal!60,rounded corners=2pt,fill=pale,align=center,font=\small,minimum height=9mm},arr/.style={-{Latex[length=2mm]},thick,draw=teal}}

\widowpenalty=10000
\clubpenalty=10000
\displaywidowpenalty=10000
\widowpenalties 3 10000 10000 0
\clubpenalties 3 10000 10000 0
\interfootnotelinepenalty=10000
\setlength{\textfloatsep}{13pt plus 2pt minus 2pt}
\setlength{\floatsep}{11pt plus 2pt minus 2pt}
\setlength{\intextsep}{12pt plus 2pt minus 2pt}
\renewcommand{\topfraction}{.9}
\renewcommand{\bottomfraction}{.8}
\renewcommand{\textfraction}{.1}
\renewcommand{\floatpagefraction}{.65}
\AddToHook{cmd/section/before}{\Needspace{7\baselineskip}}
\AddToHook{cmd/subsection/before}{\Needspace{6\baselineskip}}
\newcommand{\demopic}[4][.95]{%
  \begin{figure}[H]\centering
  \includegraphics[width=#1\linewidth,height=.56\textheight,keepaspectratio]{images/#2}
  \caption{#3}\label{#4}\end{figure}}
\newcommand{\howto}[1]{\par\smallskip\noindent\textbf{Cách sử dụng.} #1\par}
\begin{document}
\begin{titlepage}
\thispagestyle{empty}
\begin{tikzpicture}[remember picture,overlay]
\draw[vtx,line width=1.3pt] ($(current page.north west)+(13mm,-13mm)$) rectangle ($(current page.south east)+(-13mm,13mm)$);
\draw[vtx!55,line width=.35pt] ($(current page.north west)+(15mm,-15mm)$) rectangle ($(current page.south east)+(-15mm,15mm)$);
\end{tikzpicture}
\centering
\vspace*{0mm}
{\footnotesize\sffamily\href{https://github.com/Mr-Gwenchana/VDT-Phase2-Project}{github.com/Mr-Gwenchana/VDT-Phase2-Project}\par}
\vspace{4mm}
\includegraphics[width=55mm]{icons/Viettel_Aerospace_Institute_(VTX)_logo.png}\par
\vspace{7mm}
{\fontsize{14}{16}\selectfont \textbf{\color{ViettelRed}TẬP ĐOÀN CÔNG NGHIỆP -- VIỄN THÔNG QUÂN ĐỘI}}\\[0.15cm]
\vspace{3mm}
{\normalsize\bfseries VIETTEL DIGITAL TALENT 2026\par}
{\normalsize GIAI ĐOẠN 2\par}
\vspace{12mm}
{\color{vtx}\fontsize{22}{27}\selectfont\bfseries BÁO CÁO TỔNG KẾT DỰ ÁN\par}
\vspace{9mm}
{\fontsize{19}{26}\selectfont\bfseries PHẦN MỀM MÔ PHỎNG HẢI CHIẾN\par TRÊN NỀN ĐỊA HÌNH THỰC TẾ\par}
\vspace{5mm}
{\color{graytext}\normalsize Unreal Engine 5.6 / Cesium / C++ và Blueprint\par}
\vspace{10mm}
{\color{vtx}\rule{.8\linewidth}{.6pt}}\par
\vspace{7mm}
\begin{tabular}{@{}p{44mm}p{78mm}@{}}
\textbf{Thành viên thực hiện:} & Lê Việt Thành\\[2mm]
 & Nguyễn Hữu Hiệu\\[5mm]
\textbf{Mentor hướng dẫn:} & Tuấn Anh\\[2mm]
 & Ngọc Dương\\[5mm]
\textbf{Chương trình:} & Viettel Digital Talent 2026\\
\end{tabular}
\vfill
{\normalsize Hà Nội, tháng 11 năm 2026\par}
\vspace{3mm}
\end{titlepage}

\thispagestyle{empty}
\fronttitle{Ý KIẾN ĐÁNH GIÁ}
\noindent Ý kiến đánh giá:\par\vspace{3mm}
\foreach \n in {1,...,10}{\noindent\makebox[\linewidth]{\dotfill}\par\vspace{8mm}}
\vspace{5mm}
\noindent Điểm số: \dotfill\hspace{8mm}Điểm chữ: \dotfill\par
\vspace{10mm}
\begin{flushright}
Hà Nội, ngày \ldots\ tháng \ldots\ năm 20\ldots\par
\vspace{3mm}\textbf{Người đánh giá}\par(Ký, ghi rõ họ tên)
\end{flushright}
\clearpage
\thispagestyle{empty}
\fronttitle{LỜI CẢM ƠN}
Chúng em, Lê Việt Thành và Nguyễn Hữu Hiệu, xin trân trọng cảm ơn Viettel Aerospace và chương trình Viettel Digital Talent 2026 đã tạo điều kiện để nhóm thực hiện dự án \textit{Phần mềm mô phỏng hải chiến trên nền địa hình thực tế}. Đây là cơ hội để chúng em vận dụng kiến thức lập trình, đồ họa máy tính và thuật toán vào một sản phẩm tương tác có nhiều hệ thống cùng hoạt động.

Chúng em xin gửi lời cảm ơn tới mentor Tuấn Anh và mentor Ngọc Dương đã hướng dẫn và góp ý trong quá trình xây dựng sản phẩm. Những nhận xét về chuyển động thân tàu, độ dập dềnh trên sóng, thông báo tiêu diệt và giao diện kết thúc trận đã giúp nhóm hoàn thiện trải nghiệm trình diễn, đồng thời hiểu rõ hơn cách đánh giá một tính năng từ góc nhìn người sử dụng.

Quá trình thực hiện giúp nhóm tích lũy kinh nghiệm kết hợp C++ với Blueprint, tích hợp dữ liệu địa lý, điều khiển phương tiện bằng lực và tổ chức hành vi AI. Các tình huống như tàu mắc ở cửa đảo, không quay lại điểm tuần tra hoặc mất quyền ngắm bắn khi phóng to bản đồ cũng cho nhóm nhiều bài học về sự phối hợp giữa các thành phần.

Chúng em cảm ơn các anh chị tư vấn và hội đồng đánh giá đã dành thời gian xem xét sản phẩm. Nhóm mong nhận được thêm góp ý để tiếp tục hoàn thiện mô hình chuyển động, hành vi chiến thuật và khả năng mở rộng của hệ thống.
\vspace{10mm}
\begin{flushright}Nhóm thực hiện\\[3mm]\textbf{Lê Việt Thành -- Nguyễn Hữu Hiệu}\end{flushright}
\clearpage
\thispagestyle{empty}
\fronttitle{LỜI MỞ ĐẦU}
Mô phỏng hải chiến trên địa hình thực tế là bài toán kết hợp môi trường địa lý, chuyển động vật lý và hành vi tự động. Địa hình đảo quyết định nơi tàu có thể đi qua và những hướng quan sát bị che khuất. Sóng biển tác động lên thân tàu, trong khi quán tính khiến một lệnh bẻ lái không thể lập tức đổi hướng di chuyển. Vì vậy, đường đi tìm được trên bản đồ cần được chuyển thành các lệnh ga và bánh lái phù hợp với phương tiện.

Nhóm xây dựng dự án bằng Unreal Engine 5.6 và Cesium, lấy khu vực Quảng Ninh làm bối cảnh. Sản phẩm hỗ trợ hai đội tàu, cho phép người sử dụng trực tiếp điều khiển một tàu hoặc quan sát toàn bộ đội hình AI. Vị trí xuất phát và tuyến tuần tra có thể được bố trí thủ công; ngoài ra, hệ thống có thể sinh bố trí theo seed để tạo các kịch bản trình diễn khác nhau.

Các nội dung chính gồm tích hợp địa hình và mặt nước, xây dựng các loại tàu, điều khiển lực đẩy và bánh lái, mô phỏng nổi trên sóng, tìm đường A*, tránh vật cản, nhận diện mục tiêu, bắn đón, xử lý sát thương và vòng đời trận đấu. Giao diện thiết lập, HUD, minimap, thông báo tiêu diệt và bảng kết quả giúp người sử dụng theo dõi các hệ thống này trong khi chơi.

Báo cáo được tổ chức theo tính năng của sản phẩm. Mỗi nhóm chức năng trình bày mục đích, cách hoạt động, cơ sở tính toán và cách sử dụng; hình ảnh được đặt cạnh phần nội dung tương ứng. Phần cuối tổng hợp kết quả chạy thử, những giới hạn còn lại và hướng phát triển. Qua đó, người đọc có thể hiểu nhóm đã xây dựng những gì và tự thực hiện một kịch bản demo từ khâu bố trí đến khi trận đấu kết thúc.
\clearpage
\pagenumbering{roman}
\chapter*{Tóm tắt dự án}\addcontentsline{toc}{chapter}{Tóm tắt dự án}
Sản phẩm là môi trường mô phỏng hải chiến có từ một đến ba tàu mỗi đội trên địa hình Quảng Ninh. Người sử dụng lựa chọn loại tàu Small, Medium hoặc Large; thiết lập độ khó, chế độ trận và cách điều khiển; sau đó bắt đầu trận với bố trí thủ công hoặc bố trí sinh theo seed. Trong trận, tàu di chuyển bằng lực vật lý, nổi trên mặt nước và tương tác với đảo, tàu khác cùng đạn pháo.

Hệ thống AI kết hợp A* để tìm tuyến qua vùng nước với bộ điều khiển bám tuyến, giảm tốc trước khi quay, tránh vật cản và phục hồi khi không tiến triển. Chiến đấu sử dụng chung cơ chế vũ khí cho người chơi và AI: bắn đón mục tiêu, kiểm tra đường đạn, quản lý băng đạn, tính sát thương, xử lý cháy và chìm. Bộ quản lý trận thống nhất số tàu sống, điểm số và kết quả thắng thua.

Giao diện cung cấp thông tin tàu, minimap phóng đại tới 9 lần theo tàu đang quan sát, thông báo tiêu diệt và menu kết thúc trận. Sản phẩm hướng đến trình diễn kỹ thuật và thử nghiệm gameplay; các thông số vật lý và chiến đấu được hiệu chỉnh cho mô phỏng, không đại diện cho thông số khai thác của tàu quân sự thực tế.

\noindent\textbf{Mã nguồn dự án:} \url{https://github.com/Mr-Gwenchana/VDT-Phase2-Project}
\par\smallskip
\noindent\textbf{Từ khóa:} Unreal Engine, Cesium, bản đồ địa lý, tàu thủy, buoyancy, A*, AI, mô phỏng hải chiến.
\begingroup\small\setstretch{1}\tableofcontents\endgroup
\listoffigures
\listoftables
\chapter*{Danh mục ký hiệu và từ viết tắt}\addcontentsline{toc}{chapter}{Danh mục ký hiệu và từ viết tắt}
\begin{tabularx}{\linewidth}{@{}lX@{}}\toprule
\textbf{Ký hiệu}&\textbf{Ý nghĩa}\\\midrule
AI&Trí tuệ nhân tạo; trong dự án là logic điều khiển tàu tự động.\\
A*&Thuật toán tìm đường dựa trên chi phí đã đi và ước lượng tới đích.\\
Blueprint&Đồ thị lập trình trực quan và tài nguyên gameplay của Unreal Engine.\\
DEM&Mô hình số độ cao, dùng biểu diễn địa hình.\\
ECEF&Hệ tọa độ địa tâm cố định theo Trái Đất.\\
ENU&Hệ trục tiếp tuyến East--North--Up: Đông, Bắc và hướng lên.\\
HUD&Bảng thông tin hiển thị trong trận đấu.\\
LOS&Đường quan sát trực tiếp giữa hai vị trí.\\
PIE&Phiên chạy thử trong Unreal Editor.\\
HP&Điểm sức bền của tàu trong mô phỏng.\\
$\Delta t$&Khoảng thời gian giữa hai bước cập nhật, đơn vị giây.\\
$v,\omega$&Vận tốc tịnh tiến và vận tốc góc của thân tàu.\\
Roll/Pitch/Yaw&Lắc ngang, chúc/ngẩng mũi và quay hướng quanh trục đứng.\\\bottomrule
\end{tabularx}
\clearpage\pagenumbering{arabic}

\chapter{Tổng quan sản phẩm và kiến trúc hệ thống}
\section{Bài toán và mục tiêu xây dựng}
Mục tiêu của dự án là xây dựng một vòng chơi hoàn chỉnh: thiết lập chiến trường, điều khiển hoặc quan sát tàu, giao chiến và xem kết quả. Địa hình thực tế được sử dụng như một thành phần của gameplay. Một đảo có thể chia cắt tuyến đi, che khuất tầm nhìn hoặc buộc tàu phải chuyển hướng tiếp cận; kích thước thân tàu cũng ảnh hưởng đến hành lang có thể sử dụng.

Nhóm tập trung vào ba yêu cầu. Thứ nhất, các tàu phải di chuyển liên tục bằng lực, có quán tính và khả năng tiến, lùi, quay đầu. Thứ hai, AI phải dùng được cùng phương tiện với người chơi, đi được tuyến đã bố trí và tham gia chiến đấu. Thứ ba, người trình diễn phải chủ động thay đổi đội hình, điểm xuất phát, mục tiêu và chế độ trận mà không cần sửa thuật toán.

\demopic[.50]{details_ship.jpg}{Tàu người chơi trên biển Quảng Ninh: HUD ở góc trái hiển thị vị trí, tốc độ, HP, hướng và đạn; minimap ở góc phải cho biết vị trí đội hình.}{fig:overview}

\section{Các nhóm chức năng đã xây dựng}
Các tính năng được tổ chức thành sáu nhóm liên kết với nhau. Bảng dưới đây tóm tắt nội dung và vai trò của từng nhóm trong một trận mô phỏng.

\begin{table}[H]\centering\small
\caption{Các nhóm chức năng chính của sản phẩm.}
\begin{tabularx}{\linewidth}{@{}p{31mm}XX@{}}\toprule
\textbf{Nhóm}&\textbf{Nội dung}&\textbf{Giá trị sử dụng}\\\midrule
Bản đồ&Địa hình, mặt nước, tọa độ và sector.&Tạo không gian di chuyển và giao chiến.\\
Thiết lập&Đội hình, loại tàu, độ khó, seed và tuyến A/B.&Thiết kế kịch bản demo.\\
Tàu và vật lý&Lực đẩy, bánh lái, tốc độ và buoyancy.&Di chuyển có quán tính, phản ứng theo sóng.\\
AI điều hướng&A*, bám tuyến, tránh vật cản và phục hồi.&Tự đi tới mục tiêu trên nước.\\
Chiến đấu&Nhận diện, bắn đón, đạn, HP và cháy/chìm.&Tạo giao chiến và vòng đời tàu.\\
Giao diện&Thiết lập, HUD, minimap, kill feed và kết quả.&Hỗ trợ điều khiển và quan sát.\\\bottomrule
\end{tabularx}\end{table}

\section{Tổ chức C++ và Blueprint}
C++ đảm nhiệm các thuật toán và quy tắc dùng chung như tìm đường, xử lý tọa độ, giới hạn lực đẩy, kiểm tra phát bắn, quản lý HP và kết quả trận. Blueprint liên kết mesh, collider, camera, pháo, projectile, âm thanh và hiệu ứng. Cách tổ chức này cho phép hiệu chỉnh tài nguyên bằng Editor, đồng thời giữ các quy tắc quan trọng ở một nơi.

Một tàu được cấu thành từ thân vật lý, mô hình hiển thị, hệ thống nổi, chuyển động, điều hướng, vũ khí và định danh. Thân vật lý chịu lực và va chạm; mô hình hiển thị mô tả hình dạng cùng chuyển động của pháo. Hai thành phần này có liên hệ nhưng không nhất thiết dùng cùng một mesh. Ví dụ, tàu Small sử dụng hình dạng riêng trong khi tận dụng cơ cấu điều khiển pháo đã có.

\begin{figure}[H]\centering
\begin{tikzpicture}[node distance=5mm]
\node[box,text width=12cm] (a) {Thiết lập trận và dữ liệu địa lý};
\node[box,text width=5.5cm,anchor=north west,at={([yshift=-5mm]a.south west)}] (b) {AI: nhiệm vụ, mục tiêu\\Tìm đường và chiến thuật};
\node[box,text width=5.5cm,anchor=north east,at={([yshift=-5mm]a.south east)}] (c) {Người chơi\\Input, camera và ngắm bắn};
\node[box,text width=12cm,below=25mm of a] (d) {Các thành phần dùng chung của tàu\\Lực đẩy, bánh lái, buoyancy, vũ khí và HP};
\node[box,text width=12cm,below=of d] (e) {Thế giới mô phỏng: địa hình, mặt nước, projectile và va chạm};
\node[box,text width=12cm,below=of e] (f) {Quản lý trận: đội, số tàu sống, điểm và kết quả\\Giao diện hiển thị trạng thái cho người sử dụng};
\draw[arr](a)--(b);\draw[arr](a)--(c);\draw[arr](b)--(d);\draw[arr](c)--(d);\draw[arr](d)--(e);\draw[arr](e)--(f);
\end{tikzpicture}
\caption{Quan hệ giữa thiết lập, điều khiển tàu, thế giới vật lý và giao diện.}
\end{figure}

Bộ quản lý trận giữ thông tin đội, số tàu sống, lượt và kết quả. HUD chỉ đọc để hiển thị, không tự suy ra thắng thua từ số actor đang xuất hiện trên màn hình. Tương tự, AI lựa chọn ý định di chuyển nhưng lệnh cuối vẫn đi qua bộ điều hướng và cơ cấu lực của tàu.

\section{Luồng sử dụng tổng thể}
Người sử dụng mở Match Settings, lựa chọn đội hình và cách bố trí, sau đó bắt đầu trận. Hệ thống tải dữ liệu, xác nhận vị trí và chờ thân tàu ổn định trên nước. Khi chiến trường sẵn sàng, người chơi lái tàu hoặc quan sát AI. Các phát bắn hợp lệ tạo projectile; va chạm gây sát thương; tàu hết HP chuyển sang trạng thái bị tiêu diệt. Khi một đội không còn tàu sống, bảng kết quả xuất hiện và cho phép trở lại thiết lập, chơi lại hoặc thoát.

Các chương tiếp theo phân tích từng nhóm chức năng theo luồng này, bắt đầu từ bản đồ và cách tạo chiến trường.

\chapter{Bản đồ địa lý và thiết lập chiến trường}
\section{Địa hình Quảng Ninh và mặt nước}
Bản đồ sử dụng Cesium để đưa địa hình và lớp ảnh bề mặt của khu vực Quảng Ninh vào Unreal Engine. Địa hình thể hiện hình dạng đảo, bờ và độ cao; imagery cung cấp màu bề mặt. Mặt nước thuộc hệ thống Water của Unreal, tạo bề mặt biển và cung cấp thông tin độ cao nước cho các điểm nổi của tàu. Hai hệ thống có vai trò khác nhau: Cesium biểu diễn môi trường địa lý, còn Water tham gia tương tác vật lý với phương tiện.

Dự án sử dụng bộ dữ liệu \textbf{UnrealTerrainSet}, với khu vực Quảng Ninh đặt tại \path{Data/UnrealTerrainSet/QuangNinh}. Bộ dữ liệu có hai phần: \texttt{Terrain} cung cấp hình học độ cao và \texttt{Imagery} cung cấp ảnh phủ bề mặt. Cesium tải các tile theo vị trí quan sát và mức chi tiết. Màu nước nhìn thấy không tự quyết định vùng tàu đi được; điều hướng còn sử dụng dữ liệu nước riêng và kiểm tra collision của thân.

Hình dạng quần đảo tạo ra nhiều trường hợp demo: tuần tra dọc hành lang nước, tránh đảo giữa hai điểm, mất tầm nhìn khi đối phương đi sau bờ và giao cắt đội hình ở vùng nước mở. Đây là lý do bản đồ ảnh hưởng trực tiếp đến cả AI di chuyển lẫn chiến đấu.
\howto{Mở level \texttt{MainMapLevel}; chờ địa hình được tải trước khi bố trí hoặc quan sát. Đặt tàu trong vùng nước đủ rộng cho toàn bộ thân và khoảng quay, không chỉ đặt tâm actor vào phần có màu xanh.}

\section{Cấu trúc UnrealTerrainSet và luồng nạp bản đồ}
Các tệp địa hình và ảnh không được nhập thành từng asset trong Content Browser. Cesium đọc một tệp mô tả, từ đó tìm các tile con tương ứng. Cấu trúc dữ liệu sử dụng trong dự án là:
\begin{lstlisting}
Data/UnrealTerrainSet/QuangNinh/
  Terrain/
    layer.json
    {z}/{x}/{y}.terrain
  Imagery/
    tilemapresource.xml
    {z}/{x}/{y}.png
Content/SeaNavigation/
  halong_navigation.json
  quangninh_sectors.json
  quangninh_demo_routes.json
\end{lstlisting}

\textbf{Terrain/layer.json} là điểm vào của địa hình. Tệp mô tả định dạng \texttt{quantized-mesh-1.0}, hệ chiếu \texttt{EPSG:4326}, sơ đồ tile \texttt{tms}, phạm vi địa lý và những tile có sẵn ở từng mức. Trường \texttt{tiles} chứa mẫu URL \verb|{z}/{x}/{y}.terrain?v={version}|; $z$ là mức chi tiết, còn $x,y$ xác định tile. Dữ liệu hiện có các mức 0--13. Các tệp \texttt{.terrain} chứa lưới độ cao lượng tử hóa, không phải texture màu.\cite{quantized}

Một số trường chính trong tệp được trích như sau:
\begin{lstlisting}
{
  "format": "quantized-mesh-1.0",
  "scheme": "tms",
  "projection": "EPSG:4326",
  "tiles": ["{z}/{x}/{y}.terrain?v={version}"],
  "bounds": [107.260192, 20.766175, 107.576040, 21.061271]
}
\end{lstlisting}
Đây là phần trích rút gọn; tệp thật còn chứa \texttt{available}, phần mở rộng và metadata. Thứ tự \texttt{bounds} là kinh độ tây, vĩ độ nam, kinh độ đông và vĩ độ bắc. Các giá trị trên cho biết phạm vi dữ liệu, không phải tọa độ xuất phát của tàu.

\textbf{Imagery/tilemapresource.xml} mô tả lớp ảnh TMS trong hệ \texttt{EPSG:4326}. Nó khai báo tile PNG $256\times256$ pixel, phạm vi phủ và các mức 6--14. Các tệp \texttt{Imagery/\{z\}/\{x\}/\{y\}.png} là ảnh bề mặt được phủ lên hình học terrain. Vì vậy chỉ có \texttt{layer.json} mà thiếu tile con sẽ không tạo được địa hình đầy đủ; chỉ có ảnh PNG cũng không tạo được độ cao đảo.

Trong level, actor \texttt{QuangNinh\_Terrain} dùng Cesium 3D Tileset, chọn nguồn \textbf{From URL} và đọc \path{Terrain/layer.json}. Component ảnh \texttt{QuangNinh\_Imagery} dùng TMS Raster Overlay, đọc \path{Imagery/tilemapresource.xml}. Cesium dựng hình học và phủ ảnh lên bề mặt; Water cung cấp mặt biển cùng tương tác nổi riêng.\cite{cesiumfaq}

\subsection{Đường dẫn động và dữ liệu rời khi đóng gói}
Actor địa hình có component \texttt{PortableLocalTerrain}. Hai thuộc tính cấu hình lưu đường dẫn tương đối đến \path{Data/UnrealTerrainSet/QuangNinh/Terrain/layer.json} và \path{Data/UnrealTerrainSet/QuangNinh/Imagery/tilemapresource.xml}. Khi component đăng ký, hệ thống lấy thư mục project hiện hành, ghép đường dẫn rồi tạo URL \texttt{file:///} cho Cesium. Khoảng trắng và ký tự Unicode được mã hóa phù hợp trong URL.

Có thể mô tả cơ chế bằng $P_{\text{tệp}}=P_{\text{project}}+P_{\text{tương đối}}$, trong đó dấu cộng biểu diễn thao tác ghép đường dẫn. Khi chuyển project hoặc bản game sang ổ đĩa/thư mục khác, $P_{\text{project}}$ thay đổi còn cấu trúc dữ liệu tương đối được giữ. Người dùng không phải nhập lại đường dẫn ổ E:/F: của máy phát triển. Nếu thiếu metadata hoặc đường dẫn ra ngoài thư mục project, component báo lỗi thay vì tiếp tục dùng URL cũ.

Các tệp terrain/imagery được khai báo dạng \textbf{NonUFS} trong cấu hình build: chúng được phân phối thành file rời để Cesium mở qua filesystem. Ba JSON gameplay bên dưới vẫn thuộc \textbf{UFS}. Trong quá trình cook, URL preview của actor được xóa; khi bản game chạy, component dựng lại URL theo thư mục cài đặt. Cách tổ chức này giữ một nguồn UnrealTerrainSet trong Data, không yêu cầu nhân đôi bộ tile trong Content.
\howto{Giữ nguyên cây \texttt{Data/UnrealTerrainSet/QuangNinh} khi di chuyển project hoặc phân phối game. Mở MainMapLevel và chạy bình thường; component tự cấu hình terrain và imagery. Khi package, cần giữ đầy đủ thư mục kết quả cùng dữ liệu được stage, thay vì chỉ lấy một file thực thi.}

\section{Ba JSON cho điều hướng và bố trí tàu}
Các JSON trong \path{Content/SeaNavigation} phục vụ gameplay, độc lập với tệp \texttt{layer.json} dùng để vẽ địa hình. Chúng được đọc theo đường dẫn tương đối từ thư mục Content và được khai báo là dữ liệu runtime trong cấu hình build.

\textbf{halong\_navigation.json} lưu lưới nước/đất và dữ liệu khoảng trống cho A*. Tệp khai báo \texttt{VDT\_SEA\_NAV}, hệ \texttt{WGS84\_ECEF\_FIXED\_ENU\_METERS}, kích thước $400\times400$, ô 100 m và biên an toàn 125 m. Gốc địa lý của chart là $(107.296416,20.992597,0)$. Các trường \texttt{origin\_ecef}, \texttt{east\_axis} và \texttt{north\_axis} xác định phép chiếu trình bày ở phần tọa độ. Hai trường \texttt{cells\_base64} và \texttt{clearance\_base64} mã hóa dữ liệu từng ô; Base64 là cách biểu diễn byte trong JSON, không phải thuật toán tìm đường.

\textbf{quangninh\_sectors.json} mô tả lưới nhiệm vụ $25\times25$, mỗi sector 100 m, tâm EN $(14500,3500)$ và phạm vi $2500\times2500$ m. Mỗi sector có vị trí, mức nước, khoảng trống, vùng liên thông và cờ phù hợp cho từng loại tàu. Các nhóm vùng như Blue defense, Red staging và vùng giao cắt giúp chọn điểm ứng viên. Tệp còn lưu mã kiểm tra của chart; nếu thay lưới điều hướng mà giữ sector cũ, hệ thống có thể từ chối vì hai dữ liệu không còn khớp.

\textbf{quangninh\_demo\_routes.json} chứa gợi ý tuyến và vị trí: vòng tuần tra Small Blue, tuyến thân lớn, vị trí tập kết Red, hướng tiếp cận và điểm phòng thủ. Tọa độ trong các mảng là East/North theo mét, không phải longitude/latitude. Trường \texttt{demo\_seed} cung cấp seed mẫu; các gợi ý vẫn phải qua kiểm tra nước, khoảng trống và collision trước khi dùng.

Luồng dữ liệu có thể tóm tắt thành: UnrealTerrainSet tạo môi trường nhìn thấy; chart tạo vùng nước cho A*; sector/tuyến gợi ý tổ chức nhiệm vụ; bộ điều khiển đưa lệnh tới thân vật lý. Dữ liệu bờ của chart được dẫn xuất từ OpenStreetMap và không có thông tin độ sâu đáy biển. Vì thế lớp nước này không phải hải đồ bảo đảm an toàn hàng hải; kiểm tra keel/collision trong mô phỏng vẫn cần thiết.

\section{Liên kết tọa độ địa lý với tọa độ mô phỏng}
Dự án sử dụng ba cách biểu diễn vị trí. Tọa độ địa lý gồm kinh độ, vĩ độ và độ cao; tọa độ ECEF biểu diễn vị trí theo một hệ trục địa tâm; tọa độ world của Unreal phục vụ actor, camera và physics. Riêng thuật toán điều hướng làm việc trên mặt phẳng Đông--Bắc cố định theo mét. Cesium đảm nhiệm chuyển đổi giữa hệ địa lý và world.\cite{cesium}

Gọi $P$ là vị trí ECEF của tàu, $O$ là gốc bản đồ, $\mathbf e$ và $\mathbf n$ là hai vectơ đơn vị hướng Đông và Bắc. Vị trí trên mặt phẳng tìm đường được tính bởi:
\begin{equation}\label{eq:enu}
 E=(P-O)\cdot\mathbf e,\qquad N=(P-O)\cdot\mathbf n.
\end{equation}
$E,N$ có đơn vị mét. Phép trừ $P-O$ tạo độ dời so với gốc, còn tích vô hướng lấy thành phần theo từng trục. Cách làm này cho phép tính khoảng cách bằng hình học phẳng trong phạm vi bản đồ, thay vì cộng trực tiếp kinh độ và vĩ độ.

Ví dụ, tàu ở $(E,N)=(100,200)$ và đích ở $(400,600)$ có khoảng cách phẳng:
\[
 d=\sqrt{(400-100)^2+(600-200)^2}=500\ \text{m}.
\]
Khi biến đích thành vị trí world, hệ thống chuyển đổi ngược qua ECEF và Cesium, đồng thời xác định cao độ mặt nước. Unreal sử dụng centimet, vì vậy $100$ m tương ứng $10\,000$ cm trong các phép quét vật lý.

\section{Biểu diễn vùng nước và khoảng cách an toàn tới bờ}
Lưới điều hướng mô tả ô có thể đi qua và ô bị chặn. Dữ liệu cơ sở có $400\times400$ ô, kích thước ô 100 m. Ngoài cờ nước/đất, mỗi ô có thông tin khoảng trống tới vật cản. Một vùng đệm quanh bờ giúp tàu tránh chọn đường áp sát đảo ngay cả khi tâm ô vẫn nằm trên nước.

Lưới này trả lời câu hỏi \textit{có hành lang nước nối hai vị trí hay không}. Kiểm tra thân tàu trả lời câu hỏi khác: \textit{tàu có kích thước và hướng hiện tại có thể đi qua hay không}. Một tuyến hợp lệ cho Small chưa chắc đủ rộng cho Medium hoặc Large. Vì vậy, hệ thống kết hợp lưới với quét hình hộp thân, kiểm tra keel và cảm biến xung quanh.

Với bố trí thủ công, vùng collision được kiểm tra quanh tàu và các điểm A/B ở độ phân giải 25 m. Hệ thống xác nhận ba chặng: xuất phát tới A, A tới B và B tới A. Nếu marker gần mép vùng đệm, một điểm tiếp cận an toàn có thể được chọn bên trong bán kính đến đích; marker gốc vẫn được giữ. Marker nằm trên đất hoặc không có hành lang kết nối sẽ bị từ chối.
\howto{Khi gặp SETUP FAILED, kiểm tra vị trí tàu, A/B, khoảng trống quanh thân và vùng liên thông trước. Dời điểm ra vùng nước thoáng là cách xử lý đúng; tăng tốc hoặc bỏ collision không giải quyết được tuyến không hợp lệ.}

\section{Lưới sector và vai trò trong bố trí tự động}
Lưới sector là lớp phân vùng nhiệm vụ, tách biệt với lưới tìm đường. Vùng kịch bản có $25\times25$ ô, mỗi ô 100 m, tương ứng cạnh 2.500 m. Các vùng logic phục vụ khu vực phòng thủ, tuần tra, cảnh giới, tập kết và hướng tiếp cận.

\begin{figure}[H]\centering
\begin{tikzpicture}[x=.20cm,y=.20cm]
\fill[blue!8] (0,0) rectangle (12,25);
\fill[red!8] (12,0) rectangle (25,25);
\fill[teal!14] (10,7) rectangle (15,18);
\draw[step=1,gray!45,thin] (0,0) grid (25,25);
\draw[teal,thick] (0,0) rectangle (25,25);
\node[fill=white,font=\small] at (5,22) {Blue};
\node[fill=white,font=\small] at (20,22) {Red};
\node[fill=white,font=\scriptsize,align=center] at (12.5,12.5) {Vùng giao\\nhiệm vụ};
\end{tikzpicture}
\caption{Nguyên lý phân vùng trên lưới sector $25\times25$. Màu biểu diễn vùng nhiệm vụ minh họa; đất/nước được xác định bởi chart và collision riêng.}\label{fig:sector}
\end{figure}

Với kích thước ô $a$ và góc gốc $(E_0,N_0)$, tọa độ ô được xác định theo:
\begin{equation}
 i=\left\lfloor\frac{E-E_0}{a}\right\rfloor,\qquad
 j=\left\lfloor\frac{N-N_0}{a}\right\rfloor.
\end{equation}
Ví dụ, với $a=100$ m, một vị trí lệch $(250,470)$ m so với góc gốc thuộc ô $(2,4)$ nếu đánh số từ 0. Việc chia sector giúp lựa chọn khu vực ứng viên; sau đó mỗi vị trí vẫn phải được xác nhận bằng dữ liệu nước và collision. Không cần tạo 625 TargetPoint thủ công trong level.

\section{Tự bố trí vị trí xuất phát và tuyến A/B}
Chế độ \textbf{GENERATE MANUALLY (A-B)} sử dụng bố trí người dùng đã lưu trong Editor. Với mỗi AI, người dùng đặt vị trí xuất phát và gán hai TargetPoint riêng. Tàu đi từ vị trí đầu tới A, tiếp tục tới B rồi lặp A--B. Đây là chức năng phù hợp khi cần chuẩn bị trước góc quay, vùng gặp nhau hoặc hướng trình diễn.

\begin{samepage}
Các bước thiết lập:
\begin{enumerate}
\item Dừng Play, mở \texttt{MainMapLevel} trong Editor.
\item Đặt các AI tại vùng nước; chỉnh hướng mũi gần với hướng chặng đầu. Đặt \texttt{PlayerStart} cho tàu người chơi.
\item Tạo hai TargetPoint khác nhau cho mỗi AI, ví dụ \texttt{TP\_Friendly\_01\_A} và \texttt{TP\_Friendly\_01\_B}.
\item Trong \texttt{ShipEnemyNavigationComponent}, gán \texttt{PatrolActorA} và \texttt{PatrolActorB}; chọn tuần tra. Nếu chỉ muốn hai điểm, để \texttt{PatrolCircuitActors} trống.
\item Lưu level và các actor đã thay đổi. Play, chọn bố trí thủ công và bắt đầu trận.
\end{enumerate}\end{samepage}

Tên TargetPoint giúp nhận biết; tham chiếu được gán trong component mới quyết định tuyến. Việc đổi tên marker không tự gán nó cho một tàu. Tương tự, di chuyển actor trong phiên Play không thay cho việc lưu bố trí ở Editor.

Bán kính đến đích của các AI trong bố trí demo là 80 m. Hai đầu tuyến nên cách nhau hơn hai lần bán kính và có không gian ngoài điểm cuối để quay đầu. Đặt A và B quá gần có thể làm hai vùng đến đích chồng nhau, khiến chuyển chặng khó quan sát.

Trong Player + AI, người chơi tự lái; A/B chỉ điều khiển AI. Trong Full AI, slot thay người chơi cũng cần tuyến. Hai marker của slot này dùng tag \texttt{Naval.PlayerPatrol.A} và \texttt{Naval.PlayerPatrol.B}; mỗi tag cần thuộc đúng một actor, và A/B phải là hai actor khác nhau.
\howto{Muốn trình diễn tuần tra rõ ràng, đặt các đội xa nhau trước, quan sát ít nhất chu trình xuất phát--A--B--A rồi mới đưa tuyến giao cắt vào vùng chiến đấu. Không đặt nhiều điểm chờ trùng một vị trí hẹp.}

\clearpage
\section{Sinh bố trí từ seed}
Chế độ \textbf{GENERATE FROM SEED} để hệ thống lựa chọn vị trí xuất phát và mục tiêu từ các vùng nhiệm vụ. Seed là đầu vào khởi tạo bộ sinh số giả ngẫu nhiên. Có thể mô tả quan hệ bố trí bằng:
\[
 \text{bố trí}=G(\text{seed},\text{đội hình},\text{dữ liệu bản đồ},\text{cấu hình}).
\]
Cùng seed cùng các đầu vào còn lại cho phép tái lập bố trí. Điều này không có nghĩa kết quả vật lý hoặc đội thắng luôn giống nhau, vì trận còn phụ thuộc điều khiển, va chạm và diễn biến giao chiến.

Seed dương được dùng trực tiếp; giá trị 0 yêu cầu tạo seed mới. Sau khi chọn vị trí ứng viên, hệ thống kiểm tra nước, footprint, khoảng cách giữa tàu, tuyến liên thông và trạng thái ổn định nổi. Số lần thử có giới hạn để tránh vòng lặp vô hạn khi cấu hình không thể bố trí.
\howto{Chọn sinh từ seed, nhập một số nguyên không âm, chọn số tàu và loại tàu rồi Start New Match. Muốn giữ nguyên tàu đã tự đặt thì dùng chế độ thủ công; seed không di chuyển các tàu trong chế độ đó.}

\chapter{Mô hình tàu và chuyển động vật lý}
\section{Các loại tàu và cấu hình trang bị}
Dự án cung cấp ba lựa chọn Small, Medium và Large thông qua profile vật lý và trang bị. Small có mô hình thân riêng, kích thước gọn và cấu hình điều khiển phù hợp. Medium và Large dùng cùng họ mô hình gốc nhưng khác khối lượng, lực đẩy, hệ số lái và cấu hình nổi. Vì vậy ba lựa chọn gameplay không đồng nghĩa ba mô hình 3D hoàn toàn độc lập.

\demopic[.84]{switch_camera_and_map.jpg}{Quan sát một tàu Small của đội Red: HUD đổi sang BP\_ShipEnemy2 trong khi các tàu khác tiếp tục hoạt động; minimap giữ thông tin vị trí đội hình.}{fig:small}

\begin{table}[H]\centering\small
\caption{Các thông số cơ sở của ba profile tàu.}
\begin{tabularx}{\linewidth}{@{}Xrrr@{}}\toprule
\textbf{Thông số}&\textbf{Small}&\textbf{Medium}&\textbf{Large}\\\midrule
Khối lượng mô phỏng (kg)&90&150&240\\
Tốc độ tiến tối đa (m/s)&29&26&26\\
Tốc độ lùi tối đa (m/s)&7&7&7\\
Tốc độ hành trình AI (m/s)&24&20&20\\
Hệ số lực đẩy&0,7&1,0&1,1\\
Hệ số bánh lái&0,5&1,0&0,8\\
Cản tuyến tính&0,5&0,5&0,9\\
Cản góc toàn thân&0,65&0,5&0,65\\
Số phao nâng đỡ&4&4&6\\
HP theo cân bằng class&38&100&100\\
Sát thương cơ sở theo class&6&10&10\\\bottomrule
\end{tabularx}\end{table}

Các giá trị trong bảng là tham số gameplay của mô phỏng. Chúng không phải khối lượng hay tốc độ kỹ thuật của tàu ngoài đời. Profile lưu trong DataAsset và có thể được người thiết kế hiệu chỉnh. Giá trị thực khi chạy còn chịu cấu hình trận; đặc biệt chế độ Standard demo thống nhất 100 HP và sát thương cơ sở 10 cho các tàu, trong khi kích thước và chuyển động vẫn theo profile.
\howto{Trong Match Settings, chọn loại tàu ở từng slot. Slot PLAYER mặc định Medium; có thể đổi sang Small hoặc Large trước khi bắt đầu trận mới. Chọn cân bằng Standard để so hành vi chiến đấu với cùng HP, hoặc cân bằng theo class để giữ chênh lệch sức bền.}

\section{Tách thân vật lý, hình dạng hiển thị và pháo}
Thân vật lý có collider và chịu lực; mô hình hiển thị thể hiện hình dạng tàu. Cách tách này cho phép thay hình dạng ngoài mà vẫn sử dụng hệ thống chuyển động và vũ khí chung. Tuy nhiên, khi thay model phải căn lại hướng mũi, kích thước collider, vị trí camera và các điểm đặt lực; chỉ thay mesh hiển thị chưa đủ tạo một loại tàu hoạt động đúng.

Tàu Small sử dụng static mesh và texture atlas riêng. Hướng mũi của tài nguyên gốc được căn lại để khớp hướng tiến trong gameplay. Cơ cấu pháo hoạt động tận dụng bộ điều khiển yaw/pitch và socket đầu nòng của rig gốc; phần mô hình điều khiển có thể ẩn trong khi các bộ phận pháo nhìn thấy vẫn bám theo nó. Khẩu pháo trang trí thuộc mesh không tự trở thành vũ khí có thể bắn.

Ưu điểm của giải pháp là giữ được chuỗi ngắm, quay pháo, projectile, âm thanh và hiệu ứng đã có. Đồng thời, mỗi loại tàu có thể hiệu chỉnh camera và trang bị riêng mà không phải xây lại toàn bộ cơ chế chiến đấu.

\section{Lực đẩy, bánh lái và quán tính}
Tàu không được dịch chuyển trực tiếp từ waypoint này sang waypoint khác. Người chơi hoặc AI tạo yêu cầu ga và bánh lái; Blueprint chuyển yêu cầu đó thành lực và mô-men tác động lên thân. Chaos tích phân trạng thái vật lý để tạo vị trí và góc quay mới.

Ở mức mô hình cơ học, quan hệ tịnh tiến và quay được mô tả bởi:
\begin{equation}
 m\dot{\mathbf v}=\sum\mathbf F,\qquad
 \frac{d\mathbf L}{dt}=\sum\boldsymbol\tau,\qquad
 \boldsymbol\tau=\mathbf r\times\mathbf F.
\end{equation}
Trong đó $m$ là khối lượng, $\mathbf v$ là vận tốc, $\mathbf L$ là động lượng góc và $\mathbf r$ là vectơ từ tâm khối lượng tới điểm đặt lực. Một lực đặt lệch tâm có thể vừa đẩy thân vừa tạo mô-men quay. Vì có quán tính, thả ga không làm tàu dừng ngay và bẻ lái không khiến hướng vận tốc đổi tức thời.

Bánh lái và lực đẩy có hệ số riêng trong profile. Tăng lực đẩy chủ yếu thay khả năng tăng tốc; tăng hệ số lái thay mức phản ứng quay. Khối lượng, damping và kích thước thân cùng ảnh hưởng kết quả, nên hai tàu có cùng tốc độ đặt vẫn có thể tăng tốc và quay đầu khác nhau.
\howto{Dùng W để tiến, S để hãm hoặc lùi, A/D để bẻ lái. Khi đang chạy nhanh, giữ S trước để giảm vận tốc tiến; tàu chỉ chuyển sang lùi sau khi qua vùng vận tốc gần 0. Khi thử quay, chừa khoảng nước cho toàn bộ thân.}

\section{Giới hạn tốc độ và hiệu chỉnh tốc độ hành trình}
Thành phần chuyển động giảm lực đẩy khi tốc độ gần giới hạn. Với vận tốc tiến có dấu $v$, tốc độ tối đa $v_{max}$ và vùng chuyển tiếp $b$, hệ số lực tiến được tính theo:
\begin{equation}
 s=\operatorname{clamp}\left(\frac{v_{max}-v}{b},0,1\right),\qquad
 \mathbf F_{ra}=s\,k_p\mathbf F_{\text{yêu cầu}}.
\end{equation}
$k_p$ là hệ số lực đẩy của profile. Ví dụ, khi $v=v_{max}$ thì $s=0$: hệ thống ngừng bổ sung lực tiến thay vì ghi đè vận tốc hiện tại. Lùi có giới hạn riêng và được giảm lực phù hợp; lực ngược hướng khi tàu đang tiến vẫn được dùng để phanh.

Tốc độ AI mong muốn còn chịu đường đi: AI giảm tốc trước khúc cua, điểm đến hoặc vật cản. Do đó \texttt{Cruise Speed} là tốc độ hành trình yêu cầu trên đoạn thuận lợi, không phải tốc độ bắt buộc ở mọi thời điểm.
\howto{Mở DataAsset \texttt{DA\_SmallShip}, \texttt{DA\_MediumShip} hoặc \texttt{DA\_LargeShip}. Chỉnh \texttt{Forward Max Speed} để đổi giới hạn tiến, \texttt{Reverse Max Speed} để đổi giới hạn lùi, \texttt{Cruise Speed} cho AI và \texttt{Propulsion Scale} cho lực đẩy. Lưu rồi bắt đầu phiên Play mới. Khi tăng tốc, cần kiểm tra lại khoảng phanh và khả năng quay tại A/B.}

\section{Nổi trên sóng bằng các phao nâng đỡ}
Hệ thống Water dùng các phao hình cầu đặt dưới thân để xấp xỉ thể tích tương tác với nước. Khi độ cao mặt nước thay đổi, độ ngập của từng phao thay đổi, làm phân bố lực nâng và mô-men thay đổi theo. Nhờ đó thân tàu có thể nhún, lắc ngang hoặc chúc/ngẩng mũi theo sóng. Đây là mô hình nổi xấp xỉ phục vụ thời gian thực.\cite{water}

Cơ sở vật lý của lực nổi là:
\begin{equation}
 F_B=\rho gV_{\text{ngập}},\qquad
 V_{\text{chỏm}}=\pi h^2\left(r-\frac h3\right),\quad 0\leq h\leq2r.
\end{equation}
$\rho$ là mật độ nước, $g$ là gia tốc trọng trường, $V_{\text{ngập}}$ là thể tích chìm. Công thức chỏm cầu giải thích vì sao độ ngập $h$ và bán kính $r$ ảnh hưởng lực đỡ của một phao hình cầu. Trong gameplay, Water còn áp hệ số, trọng số, damping và giới hạn lực; không dùng công thức lý tưởng này như lời khẳng định lực thực luôn bằng đúng $\rho gV$.

Các phao được phân bố dọc thân và hai mạn. Small/Medium có bốn phao; Large có thêm cặp giữa để hạn chế trường hợp một hõm sóng làm thân thiếu điểm đỡ. Trọng số nâng đỡ được tính lại theo vị trí phao và tâm khối lượng, với tổng trọng số bằng 1.

Nếu tăng bán kính lên hai lần, thể tích toàn cầu tăng theo $r^3$, tức tám lần. Vì vậy, dự án phối hợp bán kính với hệ số displacement thay vì chỉ tăng kích thước phao. Nhờ đó có thể giữ tiếp xúc nước tốt hơn mà không vô tình làm tàu bị đẩy lên quá mạnh.

\section{Cản nước để tàu đầm nhưng vẫn có độ dập dềnh}
Tàu cần ổn định khi đứng yên và ngắm bắn, đồng thời vẫn phản ứng tự nhiên với sóng. Dự án bổ sung cản theo vận tốc đứng, vận tốc ngang và vận tốc góc roll/pitch. Các lực được áp trước bước physics khi thân còn sống, đang mô phỏng và có tiếp xúc nước.

Để hệ số cản ít phụ thuộc bước thời gian, hệ số hiệu dụng được tính bởi:
\begin{equation}\label{eq:damping}
 k_{eff}=\frac{1-e^{-k\Delta t}}{\Delta t},\qquad
 F_z=-m v_z k_{eff}.
\end{equation}
$k$ là mức cản theo giây, $v_z$ là vận tốc đứng. Dấu âm cho biết lực chống lại hướng nhún. Nếu chỉ xét riêng lực này trong một bước tích phân, vận tốc còn xấp xỉ $v_z e^{-k\Delta t}$. Với $k=4$ s$^{-1}$ và $\Delta t=1/60$ s, vận tốc còn khoảng 93,6\% sau một bước; lực nổi và các lực khác vẫn tiếp tục tác động.

Cản ngang được đặt tại tâm khối lượng để giảm trôi ngang mà không tạo thêm mô-men lệch hướng. Với vận tốc góc trong hệ thân $(\omega_x,\omega_y,\omega_z)$, thành phần cản góc bổ sung tương ứng $(-k_r\omega_x,-k_p\omega_y,0)$. Thành phần yaw bằng 0 để tránh làm mất khả năng rẽ. Đây là gia tốc góc cản được áp qua chế độ torque phù hợp của engine, không phải khóa góc roll/pitch.

Small dùng mức cản góc nhẹ hơn để còn chuyển động theo sóng. Các lần chạy thử ghi nhận pitch có thể đạt khoảng 2--3 độ tùy trạng thái; đây là phản ứng đo được, không phải góc dao động cưỡng bức hoặc giới hạn cứng. Medium/Large được hiệu chỉnh đầm hơn và vẫn giữ điều khiển bánh lái.
\howto{Muốn thay độ đầm, chỉnh từng nhóm nhỏ trong profile: \texttt{Heave Damping}, \texttt{Roll/Pitch Damping}, \texttt{Vertical Damping Multiplier}. Sau đó thử đứng yên, bắn, tiến/lùi và rẽ cả hai phía. Không tăng mọi hệ số inertia/damping cùng lúc vì có thể làm yaw chậm rõ rệt.}

\section{Bảo vệ khi đến gần đảo và vùng nước nông}
Hệ thống phân biệt vùng đệm điều hướng với địa hình vật lý. Vùng đệm khiến AI chọn đường xa bờ; collision giữ thân không xuyên đảo. Thành phần xử lý terrain lọc riêng hình học mặt biển/skirt đã xác nhận thuộc DEM, giữ collision của đất thật. Cách này hạn chế hiện tượng thân bị kênh lên do va vào phần địa hình không phù hợp ở mặt biển.

Với người chơi, tính năng ngăn mắc cạn quét hình hộp thân theo hướng lực đẩy. Khoảng nhìn trước tăng theo tốc độ:
\begin{equation}
 H=\operatorname{clamp}(4v_++20,20,180)\quad(\text{m}),
\end{equation}
trong đó $v_+=\max(v_{\text{hướng chạy}},0)$ có đơn vị m/s; hệ số 4 biểu diễn khoảng dự báo 4 giây. Ví dụ, ở 10 m/s, khoảng quét là 60 m. Khi sắp chạm địa hình, lực tiến có thể chuyển sang hãm và HUD hiện cảnh báo nước nông.

Tính năng vẫn để người chơi rẽ hoặc lùi ra và không tự lái tàu. Nó cũng không thay thế việc đặt xuất phát hợp lệ: nếu thân đã nằm sâu trên đất, một lớp chặn ga không thể bảo đảm tự đưa tàu về nước.

\begingroup\setstretch{1.1}\setlength{\parskip}{3pt}
\chapter{AI điều hướng và hành vi chiến thuật}
\section{Tìm đường A* trên vùng nước}
A* tìm một dãy ô nước từ vị trí hiện tại tới mục tiêu. Mỗi nút có chi phí đã đi $g$ và ước lượng tới đích $h$; nút có tổng $f=g+h$ nhỏ hơn được ưu tiên mở rộng. Lưới cho phép tám hướng, gồm bốn hướng thẳng và bốn đường chéo.

Trong dự án, chi phí đi từ ô $u$ sang ô $v$ là:
\begin{align}
 c(u,v)&=d(u,v)\left(1+\frac{0{,}15}{1+C(v)}\right),\\
 f(v)&=g(v)+\|p_v-p_{\text{đích}}\|_2.
\end{align}
$d$ bằng cạnh ô $a$ khi đi thẳng, hoặc $a\sqrt2$ khi đi chéo. $C(v)$ là giá trị clearance lưu tại ô; nó tạo phạt nhẹ cho đường sát vật cản. Số hạng này là chi phí chọn tuyến của gameplay, không phải công thức tiêu hao nhiên liệu.

Ví dụ, với $a=100$ m, một bước thẳng ở ô có $C=0$ tốn 115 đơn vị chi phí; ở ô có $C=4$ tốn 103. Nếu hai tuyến gần tương đương về chiều dài, tuyến có khoảng trống tốt hơn có xu hướng được ưu tiên. Khi đi chéo, cả hai ô cạnh liên quan cũng phải hợp lệ để tránh lách qua góc đất.

Sau khi tới đích, thuật toán truy vết nút cha để khôi phục tuyến. Bước rút gọn thử nối thẳng các điểm xa hơn nhưng vẫn kiểm tra cùng lưới nước và các vùng loại trừ; nó không được phép cắt xuyên bờ. Kết quả là danh sách waypoint gọn hơn để bộ điều khiển bám theo.

Mỗi lần giải có giới hạn số nút và thời gian: cấu hình chuẩn dùng tối đa 75.000 lần mở rộng, khoảng ngân sách 20 ms được kiểm tra theo nhóm nút, cùng giới hạn bộ nhớ open set. Khi không có tuyến nước nối thông, hệ thống báo thất bại thay vì trả đường xuyên đất.

\section{Một luồng lệnh từ nhiệm vụ tới chuyển động}
Tuần tra, chiến đấu, tìm lại địch và rút lui là các nguồn \textit{ý định}. Bộ điều hướng chuyển ý định thành mục tiêu và tuyến; cảm biến, phanh cùng tránh vật cản hiệu chỉnh lệnh; sau đó chỉ một cặp ga/bánh lái cuối được đưa vào cơ cấu lực. Nhờ vậy, hai hệ thống không đồng thời ghi đè lực đẩy theo hai mục tiêu khác nhau.

\begin{figure}[H]\centering
\begin{tikzpicture}[node distance=4mm]
\node[box,text width=11.5cm] (a) {Ý định: tuần tra A/B hoặc nhiệm vụ chiến thuật};
\node[box,text width=11.5cm,below=of a] (b) {A*: tuyến qua vùng nước và các waypoint};
\node[box,text width=11.5cm,below=of b] (c) {Bám tuyến, giảm tốc, cảm biến và tránh giao thông};
\node[box,text width=11.5cm,below=of c] (d) {Lệnh cuối: throttle và rudder};
\node[box,text width=11.5cm,below=of d] (e) {Blueprint áp lực và mô-men\\Chaos cùng buoyancy cập nhật thân tàu};
\draw[arr](a)--(b);\draw[arr](b)--(c);\draw[arr](c)--(d);\draw[arr](d)--(e);
\end{tikzpicture}
\caption{Luồng điều khiển AI thống nhất; tuyến hình học được chuyển thành tác động vật lý.}
\end{figure}

A* không được gọi vô điều kiện mỗi frame. Hệ thống giữ tuyến đang dùng và yêu cầu lập lại khi đích thay đổi, đoạn phía trước không còn hợp lệ hoặc quá trình phục hồi cần đường mới. Điều này giảm công việc lặp lại và tránh việc tàu liên tục đổi ý định trên một tuyến vẫn sử dụng được.

\section{Bám waypoint, giảm tốc và quay đầu}
Bộ điều khiển xác định sai lệch góc giữa hướng mũi $\mathbf f$ và hướng tới waypoint $\mathbf d$ trên mặt phẳng thân. Với $\mathbf u$ là hướng lên của thân, góc có dấu được tính bằng:
\begin{equation}
 e_\psi=\operatorname{atan2}\bigl(\mathbf u\cdot(\mathbf f\times\mathbf d),\mathbf f\cdot\mathbf d\bigr).
\end{equation}
Góc được đổi sang độ để phát lệnh bánh lái. Dạng điều khiển thực hiện kết hợp sai lệch góc và cản tốc độ quay:
\[
 r=\operatorname{clamp}\left(s_r\frac{e_\psi}{45}-k_\omega\omega_{yaw},-1,1\right).
\]
$s_r$ điều chỉnh chiều bánh lái, $k_\omega$ giảm quay quá mức. Khi sai lệch khoảng 45 độ, thành phần theo góc đã đạt mức lái lớn; khi thân đang quay nhanh, phần cản giúp hạn chế vượt hướng.

Tàu cần giảm tốc trước khi tới waypoint cuối. Từ quan hệ phanh chuyển động thẳng $v^2=2ad$, hệ thống giới hạn tốc độ theo:
\begin{equation}
 v_{\text{đích}}\leq\sqrt{2a_b\max(0,d_{\text{còn}}-R_{\text{đến}})}.
\end{equation}
$a_b$ là giảm tốc phanh cấu hình, $d_{\text{còn}}$ là độ dài tuyến còn lại và $R_{\text{đến}}$ là bán kính chấp nhận. Ví dụ còn 100 m ngoài vùng đến với $a_b=1$ m/s$^2$, tốc độ theo điều kiện phanh không nên vượt khoảng 14,14 m/s. Ở góc cua, AI còn xét bán kính quay và gia tốc ngang; quan hệ $v_{cua}\approx\sqrt{a_{ngang}R}$ giải thích vì sao tàu phải chậm hơn khi đổi hướng gấp.

Điểm nhìn trước trên tuyến tăng theo tốc độ để tránh bám cứng từng ô nhỏ. Tuy nhiên, chỉ được bỏ qua waypoint nếu đoạn nối tới điểm tiếp theo vẫn thông thoáng. Tàu đến đúng vùng đích, hãm và chờ rồi mới chuyển chặng, giúp tránh dao động quanh một điểm ở tốc độ cao.

\section{Tuần tra A--B lặp lại}
Trạng thái tuần tra lưu điểm đang hoạt động, số chặng hoàn tất và thời gian chờ. Chu trình mong muốn là:
\[
 S\longrightarrow A\longrightarrow B\longrightarrow A\longrightarrow B\longrightarrow\cdots
\]
$S$ là vị trí xuất phát đã đặt. Nếu $S$ nằm sẵn trong bán kính A, chặng đầu có thể kết thúc nhanh. Đây là lý do nên bố trí khoảng xuất phát rõ ràng khi trình diễn.

\begin{figure}[H]\centering
\begin{tikzpicture}[node distance=18mm]
\node[box,text width=25mm] (s) {Vị trí đầu};
\node[box,text width=25mm,right=of s] (a) {Target A};
\node[box,text width=25mm,right=of a] (b) {Target B};
\draw[arr] (s)--node[above,font=\small]{chặng đầu}(a);
\draw[arr] (a) to[bend left=25] node[above,font=\small]{đi tới B} (b);
\draw[arr] (b) to[bend left=25] node[below,font=\small]{quay lại A} (a);
\end{tikzpicture}
\caption{Chu trình tuần tra: từ vị trí xuất phát tới A, sau đó lặp A--B. Mỗi chặng được lập đường và thực hiện bằng lực vật lý.}\label{fig:routes}
\end{figure}

Trong chế độ thủ công, A/B là tuyến do người dùng thiết kế. AI vẫn phát hiện và bắn đối phương, nhưng không tự thay tuyến bằng truy đuổi hoặc rút lui chiến thuật. Nhờ đó, người trình diễn có thể chuẩn bị đường đi và thời điểm gặp nhau. Trong bố trí sinh tự động, quyền thay đổi ý định chiến thuật được mở để AI phản ứng với tình huống.
\howto{Theo dõi tàu thực sự tới cả A và B; không chỉ nhìn một lần quay mũi. Nếu tàu dừng bất thường, kiểm tra cặp tham chiếu, khoảng cách hai điểm, bán kính đến và không gian quay trước khi thay thuật toán.}

\section{Tránh đảo, tránh tàu và quyền ưu tiên giao thông}
Đường A* mô tả vật cản tĩnh, còn tàu khác là vật cản động. Bộ điều khiển dùng các cảm biến hướng trước và hai bên, kết hợp quét thân để phát hiện nguy cơ. Lệnh cuối có thể giảm ga, hãm, chọn phía tránh hoặc yêu cầu một tuyến vòng.

Khoảng cần thiết để dừng được ước lượng theo vận tốc:
\begin{equation}
 d_{\text{dừng}}=\frac{v^2}{2a_b}+2v+m_s.
\end{equation}
Số hạng đầu là quãng phanh lý tưởng, $2v$ bổ sung khoảng chạy trong hai giây và $m_s$ là biên an toàn. Với $v=10$ m/s, $a_b=1$ m/s$^2$ và $m_s=10$ m, khoảng dự báo là 80 m. Tăng tốc gấp đôi làm phần quãng phanh tăng gấp bốn, nên tăng Cruise Speed cần đi cùng kiểm tra hành lang.

Khi gặp tàu cùng hướng, hệ thống phân biệt tàu đang dẫn đoàn với tàu đứng yên. Trường hợp thứ nhất phù hợp giảm tốc bám theo; trường hợp thứ hai có thể cần chờ hoặc tìm đường vượt. Cách phân biệt này hạn chế việc một đoàn AI đứng chờ vô thời hạn sau một tàu không di chuyển.

Vùng tàu cản có thể được đưa vào planner như vùng loại trừ tạm. Khi tình huống thay đổi, vùng loại trừ có thời hạn và được đánh giá lại. Hệ thống vẫn giữ kiểm tra nước và footprint; tránh tàu không được làm đường mới xuyên qua đảo.

\section{Phục hồi khi không còn tiến triển}
AI theo dõi việc giảm khoảng cách còn lại trên tuyến. Nếu tàu dùng ga nhưng không tiến triển đủ trong cửa sổ quan sát, trạng thái có thể chuyển sang tránh, hãm, lùi kiểm tra hoặc lập lại đường. Hướng lùi cũng cần được xác nhận, không phải cứ mắc là lùi thẳng bất kể vật cản phía sau.

Những hành lang đã thử thất bại được ghi nhớ tạm để tránh lập đi lập lại cùng một tuyến xấu. Sau thời gian chờ, hệ thống có thể thử lại vì tàu khác đã rời đi hoặc điều kiện collision đã thay đổi. Số lần thử và thời gian chờ giúp tránh lặp dày trong một frame, đồng thời không biến một thất bại tạm thời thành trạng thái đứng yên vĩnh viễn.


\section{Nhiệm vụ chiến thuật và tìm lại mục tiêu}
Trong bố trí tự động, Blue có nhiệm vụ phòng thủ hoặc tuần tra; Red có thể tiếp cận theo nhóm với tàu dẫn và tàu theo. Nhiệm vụ chọn khu vực cần tới, còn bộ điều hướng quyết định đường đi và lệnh ga/bánh lái.

Vai trò chiến đấu quyết định mức chủ động: Passive hạn chế tấn công theo cấu hình; Coward ưu tiên tách nguy cơ; Fighter giao chiến trong phạm vi; Hunter truy đuổi có giới hạn thời gian, khoảng cách và HP. Các vai trò này khác nhiệm vụ đội hình Scout hay Follower.

Khi địch bị đảo che, AI lưu vị trí, vận tốc đã quan sát và thời điểm. Tìm kiếm dùng ký ức còn hiệu lực cùng các điểm nước hợp lệ, thay vì biết vị trí mới xuyên địa hình. Tái thu nhận được địch thì tiếp tục ngắm/bắn; tìm kiếm không tiến triển thì trở về nhiệm vụ.

Nếu được phép thay tuyến, AI có thể rút khi HP thấp hoặc địch quá gần. Nó kiểm tra phía sau trước khi lùi, tìm đường thoát tiến nếu bị chặn và trở lại nhiệm vụ khi an toàn. Chế độ thủ công giữ A/B nên không tự bỏ tuyến để truy đuổi/rút lui.

\par\endgroup
\chapter{Vũ khí, sát thương và vòng đời trận đấu}
\section{Phát hiện mục tiêu và kiểm tra tầm nhìn}
AI lựa chọn mục tiêu từ các tàu còn sống, khác đội và trong phạm vi hoạt động. Tầm phát hiện dùng để nhận biết đối phương, còn tầm giao chiến giới hạn việc bắn. Cấu hình demo sử dụng các mức khoảng cách như phát hiện 1.800 m, giao chiến 1.200 m và mất mục tiêu 2.200 m; các giá trị này có thể hiệu chỉnh theo kịch bản.

Đối phương nằm trong bán kính chưa chắc bắn được. LOS kiểm tra đường quan sát có bị đảo hoặc địa hình chắn hay không. Sau đó hệ thống còn kiểm tra góc pháo, nghiệm đạn đạo, đường bay và quyền bắn. Nhờ vậy, khoảng cách và hình dạng địa hình cùng tham gia quyết định một phát bắn.
\howto{Đặt hai tàu khác đội tiến vào vùng nước mở để quan sát bắt mục tiêu; sau đó dùng một đảo ở giữa để quan sát mất LOS và tìm lại địch. Không dùng hai tàu đồng đội để thử sát thương đối phương.}

\demopic[.86]{fighting.jpg}{Góc nhìn ngắm pháo khi hai đội ở gần nhau. Tâm ngắm, thanh HP và nhãn đội cho phép theo dõi mục tiêu trong lúc giao chiến.}{fig:aimdemo}

\section{Bắn đón mục tiêu chuyển động}
Nếu nhắm đúng vị trí hiện tại của mục tiêu, đạn có thể tới nơi sau khi tàu đã rời đi. Bộ giải bắn đón ước lượng thời gian bay $T$, dự đoán vị trí tương lai của địch rồi tìm vận tốc phóng tương ứng.

Gọi $\mathbf p_0$ là đầu nòng, $\mathbf p_t$ là vị trí quan sát của mục tiêu, $\mathbf v_t$ là vận tốc mục tiêu và $\mathbf g$ là gia tốc trọng trường. Mô hình dùng trong bộ giải là:
\begin{align}
 \mathbf p_{\text{đón}}(T)&=\mathbf p_t+\mathbf v_tT,\\
 \mathbf p_{\text{đạn}}(T)&=\mathbf p_0+\mathbf v_0T+\tfrac12\mathbf gT^2.
\end{align}
Trong một khoảng dự báo ngắn, mục tiêu được xấp xỉ chuyển động đều. Thuật toán lặp cập nhật thời gian bay và điểm đón, gọi bộ giải vận tốc projectile của Unreal rồi kiểm tra hội tụ. Nó có giới hạn mười vòng lặp; biến thiên thời gian dưới 0,001 s được xem là đủ nhỏ, cùng kiểm tra sai số phù hợp.

Ví dụ minh họa: nếu đạn bay khoảng 2 s và mục tiêu chạy ngang 10 m/s thì điểm đón cần lệch khoảng 20 m theo hướng chuyển động, trước khi xét trọng lực và sai số. AI sử dụng lịch sử quan sát có trễ theo độ khó; vị trí và vận tốc dự báo phải thuộc cùng mẫu quan sát.

Nghiệm toán học chưa đủ để chấp nhận phát bắn. Quỹ đạo được chia thành các đoạn để kiểm tra vật cản, chính thân tàu và nguy cơ với đồng đội. Cơ cấu pháo phải quay tới đúng hướng trong giới hạn cơ khí. Projectile không tự được cộng vận tốc tàu bắn nếu cơ chế projectile hiện có không quy định việc kế thừa đó.

\section{Một cơ chế vũ khí chung cho người chơi và AI}
Người chơi và AI dùng chung thành phần quản lý phát bắn. Yêu cầu từ chuột hoặc từ AI đều qua các điều kiện: tàu còn sống, có quyền hành động, không nạp đạn, còn đạn và không bị cooldown/recoil chặn. Khi đủ điều kiện, Blueprint tạo projectile cùng hiệu ứng và âm thanh; chỉ spawn thành công mới ghi nhận phát bắn và trừ một viên.

\begin{figure}[H]\centering
\begin{tikzpicture}[node distance=4mm]
\node[box,text width=11.7cm] (a) {Yêu cầu bắn từ người chơi hoặc AI};
\node[box,text width=11.7cm,below=of a] (b) {Kiểm tra sống, lượt, băng đạn, nạp, cooldown và hướng pháo};
\node[box,text width=11.7cm,below=of b] (c) {Blueprint sinh projectile, âm thanh và hiệu ứng};
\node[box,text width=11.7cm,below=of c] (d) {Xác nhận sinh thành công\\Trừ đạn và bắt đầu cửa sổ hồi vũ khí};
\draw[arr](a)--(b);\draw[arr](b)--(c);\draw[arr](c)--(d);
\end{tikzpicture}
\caption{Chuỗi xử lý phát bắn dùng chung giúp giữ nhất quán giữa người chơi và AI.}
\end{figure}

Cấu trúc này giúp tránh việc người chơi có một hệ đạn còn AI dùng một hệ sát thương tách biệt. Các callback lặp không được tính thành nhiều phát bắn. Projectile lưu thông tin đội và nguồn gây sát thương, nên vẫn xác định đúng quan hệ địch/ta sau khi actor bắn đã bị dọn.

\demopic[.86]{effect_fighting.jpg}{Khói từ khu vực pháo trong một tình huống khai hỏa; mô hình pháo và hiệu ứng được giữ cùng chuỗi vũ khí dùng chung.}{fig:effects}

\section{Băng đạn, nạp và cảnh báo ít đạn}
Mỗi tàu bắt đầu với băng 10 viên. Mỗi phát thành công giảm một viên; hết băng tự chuyển sang nạp trong 30 giây thời gian mô phỏng. Trong thời gian nạp, yêu cầu bắn bị chặn. Tàu chết sẽ hủy quá trình nạp; đổi mục tiêu không nạp lại băng.

Nếu $N$ là số viên còn lại, quy tắc cập nhật cơ bản là:
\[
 N_{\text{mới}}=N_{\text{cũ}}-1\quad\text{khi phát bắn được tạo thành công}.
\]
Khi hết thời gian nạp, $N$ được đưa về dung lượng băng. HUD hiển thị \texttt{AMMO: N/10} và thời gian \texttt{RELOADING}. Khi còn từ 0 đến 3 viên, dòng đạn chuyển đỏ để cảnh báo; khi nạp đầy trở lại màu xanh. Màu này chỉ biểu diễn trạng thái đạn, không thay đổi cơ chế vũ khí.

Cooldown và nạp là hai điều kiện khác nhau. Cooldown giới hạn nhịp từng phát; reload bổ sung băng sau khi hết đạn. Ngay cả khi cooldown native của người chơi bằng 0, cơ cấu recoil hoặc điều kiện bắn trong Blueprint vẫn có thể chặn yêu cầu quá sớm.
\howto{Theo dõi AMMO trước khi đưa tàu vào vùng giao chiến. Nếu không bắn được, kiểm tra nạp/cooldown, lượt và việc camera ngắm đang hoạt động. P tạm dừng thời gian mô phỏng nên đồng hồ nạp cũng dừng.}

\section{Độ khó và sai số ngắm bắn}
Độ khó tác động tới độ phân tán hướng bắn, độ trễ theo dõi, nhịp cập nhật aim và thời gian thu nhận mục tiêu. Mục tiêu là tạo mức phản ứng khác nhau, thay vì chỉ thay HP đối phương.
\begin{table}[H]\centering\small
\caption{Một số tham số của ba mức độ khó.}
\begin{tabular}{@{}lrrr@{}}\toprule
\textbf{Mức}&\textbf{Accuracy}&\textbf{Sai số aim}&\textbf{Trễ theo dõi}\\\midrule
Easy&0,35&3,0 độ&0,65 s\\
Normal&0,60&2,2 độ&0,35 s\\
Hard&0,90&0,35 độ&0,08 s\\\bottomrule
\end{tabular}\end{table}

Accuracy là tham số dùng trong mô hình phân tán, không có nghĩa xác suất trúng đo được là 35\%, 60\% hay 90\%. Hiệu quả thực tế còn phụ thuộc khoảng cách, tốc độ mục tiêu, địa hình che khuất và chuyển động thân tàu. Khi so sánh các mức, cần giữ cùng bố trí và loại tàu để tránh nhầm khác biệt của kịch bản với khác biệt của độ khó.

\section{Sát thương trực tiếp, sát thương vùng và chặn trùng}
Điểm trúng được chuyển về không gian thân để xác định vùng mũi, giữa hoặc đuôi. Sát thương trực tiếp có thể viết dưới dạng:
\begin{equation}
 D_{\text{trúng}}=\operatorname{round}(D_{\text{cơ sở}}\,k_{\text{vùng}}),\qquad
 k_{\text{vùng}}\in\{0{,}8;1;0{,}8\}.
\end{equation}
Với sát thương cơ sở 10, ba vùng tương ứng 8, 10 và 8 HP. Với Small theo cân bằng class có sát thương cơ sở 6, kết quả làm tròn là 5, 6 và 5 HP. HP mới bị giới hạn không nhỏ hơn 0.

Sát thương vùng kiểm tra nạn nhân trong bán kính và điều kiện che khuất. Nó dùng sát thương cơ sở theo nhánh riêng, không tự áp lại hệ số vùng thân như một phát trực tiếp thứ hai. Projectile có cơ chế ngăn ghi nhận nhiều hit cho cùng nạn nhân. Kiểm tra cuối cũng chặn self damage và friendly damage, giúp tránh điểm số sai do đạn chạm chính tàu hoặc tàu cùng đội.

Điểm sát thương được tính theo lượng HP thực mất:
\[
 D_{\text{ghi điểm}}=HP_{\text{trước}}-HP_{sau}.
\]
Ví dụ, tàu còn 3 HP nhận hit 10 thì chỉ ghi thêm 3 sát thương, không ghi 10 điểm overkill. Cách tính này giữ bảng kết quả phản ánh sức bền thực sự bị loại bỏ.

\Needspace{12\baselineskip}\section{Tiêu diệt, cháy, chìm và dọn xác tàu}
Tàu bị tiêu diệt khi HP về 0 hoặc thỏa điều kiện lật quá mức trong thời gian xác nhận. Sự kiện chết chỉ được tiếp nhận một lần: vũ khí dừng bắn, AI ngừng điều khiển chủ động, bộ quản lý trận giảm số tàu sống và giao diện nhận thông báo.

Chuỗi trình diễn gồm:
\[
 \text{còn sống}\rightarrow\text{bị tiêu diệt}\rightarrow\text{cháy}
 \rightarrow\text{chìm}\rightarrow\text{ngập hoàn toàn}\rightarrow\text{dọn xác}.
\]
Cấu hình cơ sở có khoảng cháy 4 s, chìm 12 s và trễ dọn 7 s. Đây là thời gian theo trạng thái, có thể thay đổi theo tài nguyên và nhịp cập nhật. Hiệu ứng cháy/nổ, âm thanh và chuyển động chìm giúp người xem nhận ra tàu đã bị loại thay vì chỉ biến mất lập tức.

Vị trí xác được theo dõi trong khung địa lý phù hợp, tránh việc thay gốc world khiến chuyển động chìm bị lệch. Khi actor được dọn, camera và HUD phải bỏ tham chiếu không còn hợp lệ. Nếu người chơi chết, hệ thống chuyển sang theo dõi đồng đội còn sống hoặc góc nhìn toàn cảnh.

\section{Thông báo tiêu diệt và xác định thắng thua}
Kill feed hiển thị tên tàu tiêu diệt và tàu bị tiêu diệt. Tên hiển thị lấy từ định danh dễ đọc, giúp người xem hiểu diễn biến ngay cả khi không nhìn trực tiếp vị trí trúng đạn. Hiệu ứng va chạm phục vụ gameplay; hình cầu debug đỏ không thuộc phần trình diễn thông thường.

\demopic[.90]{defeat_enemy.jpg}{Player Ship tiêu diệt BP\_ShipEnemy2: thông báo xuất hiện ở góc trên trái và tàu địch cháy trên mặt nước. HUD đồng thời cho thấy băng hết đạn, đang nạp và dòng AMMO màu đỏ.}{fig:kill}

Bộ quản lý trận giữ tập tàu đã chết để không giảm bộ đếm lặp. Khi một bên không còn tàu sống, kết quả được phát một lần. Player + AI hiển thị thắng/thua theo đội người chơi; Full AI cho biết Blue hoặc Red thắng. Bảng tổng kết liệt kê sát thương gây ra, sát thương nhận và số lần tiêu diệt của từng tàu.



\chapter{Giao diện và hướng dẫn sử dụng}
\section{Match Settings và lựa chọn đội hình}
Match Settings gom các quyết định trước trận vào một màn hình: độ khó, chế độ trận, cách bố trí, quyền điều khiển, đội người chơi, cân bằng và từng slot tàu. Mỗi đội có từ một đến ba tàu. Trong Player + AI, slot đầu của đội được chọn là người chơi, các slot còn lại dùng AI.

\demopic[.56]{start_menu.jpg}{Match Settings: chọn độ khó, chế độ trận, cách bố trí, quyền điều khiển và loại tàu cho hai đội. Ảnh minh họa đội hình ba Medium mỗi bên.}{fig:settings}

\begin{table}[H]\centering\small
\caption{Ý nghĩa các nhóm lựa chọn trước trận.}
\begin{tabularx}{\linewidth}{@{}p{37mm}X@{}}\toprule
\textbf{Lựa chọn}&\textbf{Cách dùng}\\\midrule
AI Difficulty&Easy, Normal hoặc Hard; thay phản ứng và sai số ngắm của AI.\\
Battle Mode&Real-Time hoặc Turn-Based với cửa sổ lượt 10 giây.\\
Deploy Troops&Chọn bố trí thủ công A/B hoặc sinh từ seed.\\
Control Mode&Player + AI để tự lái; Full AI / spectator để quan sát.\\
Player Team&Chọn đội người chơi trong chế độ có người lái.\\
Ship Mode&Standard 100 HP/10 damage hoặc cân bằng theo class.\\
Blue/Red slots&Thêm/bớt số tàu và chọn Small/Medium/Large cho từng slot.\\
Start New Match&Áp các lựa chọn vào một trận mới.\\
Quit / Back to Game&Thoát ở màn hình khởi đầu hoặc trở lại trận đang xem, tùy ngữ cảnh.\\\bottomrule
\end{tabularx}\end{table}

Sau khi Start, màn hình chuẩn bị được giữ trong lúc dữ liệu và vị trí tàu được xác nhận. Người dùng cần chờ giai đoạn này hoàn tất trước khi đánh giá việc tàu có chạy hay không. Khi thiết lập không hợp lệ, thông báo lỗi hướng người dùng trở về chỉnh bố trí thay vì tiếp tục một trận thiếu tàu.

\demopic[.86]{preparing_battlefield.jpg}{Màn hình PREPARING BATTLEFIELD trong lúc khởi tạo môi trường và vị trí tàu. Người dùng chờ bước này hoàn tất trước khi bắt đầu điều khiển.}{fig:preparing}

\clearpage\section{Điều khiển tàu và thao tác bắn}
Input được tổ chức theo các hành động lái, nhìn, ngắm và bắn. Enhanced Input giúp tách phím/chuột cụ thể khỏi hành động gameplay; controller bổ sung các phím quản lý camera và giao diện.\cite{input}

\begin{longtable}{@{}p{30mm}p{115mm}@{}}
\caption{Các thao tác điều khiển và quan sát.}\\
\toprule\textbf{Phím/thao tác}&\textbf{Chức năng}\\\midrule\endfirsthead
\toprule\textbf{Phím/thao tác}&\textbf{Chức năng}\\\midrule\endhead
W / S&Tiến; hãm hoặc lùi theo trạng thái vận tốc.\\
A / D&Bẻ lái trái/phải.\\
Di chuột&Đổi góc nhìn và hướng ngắm theo camera.\\
Giữ chuột phải&Vào camera ngắm của pháo; thả để trở lại góc thường.\\
Chuột trái&Yêu cầu bắn khi camera ngắm và vũ khí đủ điều kiện.\\
Cuộn chuột&Thay mức zoom camera khi chơi; trong tương tác minimap thì đổi mức zoom bản đồ.\\
M&Bật/tắt phóng đại minimap 1 lần/9 lần, theo tàu đang xem.\\
Giữ Left Alt&Hiện con trỏ để kéo, cuộn và thao tác trên minimap; thả để trở về gameplay.\\
Q&Chuyển tàu đang quan sát.\\
F&Chuyển góc nhìn toàn cảnh hoặc trở về tàu.\\
I&Bật/tắt bảng thông tin tàu.\\
P&Tạm dừng và mở thiết lập trận.\\
Tab&Chuyển đồng đội còn sống khi người chơi chết; chuyển tàu quan sát trong Full AI.\\
/&Mở/đóng hướng dẫn điều khiển.\\
Shift + F1&Thả chuột ra Unreal Editor trong phiên PIE.\\\bottomrule
\end{longtable}

Để thử bắn, chọn Player + AI và bắt đầu trận, đưa camera hướng ra vùng nước, giữ chuột phải để vào chế độ ngắm, rồi nhấn chuột trái khi tâm ngắm xuất hiện. Đồ thị ngắm có thể ẩn tâm nếu ray chạm chính thân tàu; khi đó cần đổi góc ra ngoài boong. Giữ nút bắn từ lúc chưa đủ điều kiện không luôn thay cho một lần nhấn mới sau đó.

Thao tác Q chỉ đổi tàu được theo dõi, không chiếm quyền lái từ AI. Trong Player + AI, W/A/S/D vẫn điều khiển tàu của người chơi. Cách tách này giữ đội hình AI tiếp tục hoạt động khi người xem đổi camera.

\section{HUD: vị trí, tốc độ, HP, hướng và đạn}
Bảng thông tin tàu hiển thị định danh, longitude, latitude, tốc độ, HP, hướng mũi, băng đạn và trạng thái nạp. Các giá trị lấy từ tàu đang quan sát, giúp đối chiếu hình ảnh với trạng thái thực. Bộ đếm hai đội lấy từ bộ quản lý trận thay vì đếm các mô hình còn nhìn thấy, vì xác tàu vẫn có thể tồn tại sau khi bị loại.

Tốc độ được đổi từ vận tốc ngang sang knots:
\begin{equation}
 v_{knots}=\frac{\sqrt{v_x^2+v_y^2}}{100}\times1{,}94384,
\end{equation}
với $v_x,v_y$ theo cm/s. Ví dụ 1.000 cm/s tương đương 10 m/s, tức khoảng 19,44 knots. Không dùng vận tốc đứng trong phép tính để sóng không bị hiểu nhầm thành tốc độ chạy tàu.

HP được thể hiện bằng số và thanh sức bền. Màu thay đổi theo mức HP giúp nhận ra tình trạng nguy hiểm; màu đội hỗ trợ phân biệt địch/ta. AMMO chuyển đỏ khi còn tối đa ba viên và kèm đồng hồ khi nạp. Cảnh báo nước nông xuất hiện khi lớp bảo vệ phía trước phát hiện nguy cơ mắc cạn.
\howto{Nhấn I để ẩn/hiện bảng khi cần quan sát rộng hơn. Khi đổi Q sang tàu khác, đọc lại tên và HP trên bảng để tránh nhầm trạng thái của tàu trước với tàu đang xem.}

\section{Minimap phóng đại và tự theo tàu}
Minimap biểu diễn vị trí tàu trên nền địa hình, dùng màu và tên để phân biệt các đối tượng. Phím M chuyển từ phạm vi tổng thể sang mức phóng đại 9 lần và tự căn theo tàu đang quan sát. Khi chuyển camera sang tàu khác, tâm theo dõi cập nhật theo marker mới; nếu không có marker được chọn phù hợp, hệ thống có thể dùng tâm của các tàu còn sống.

Gọi $\mathbf q$ là vị trí marker trong tọa độ bản đồ, $z$ là hệ số zoom, $\mathbf p$ là độ pan và $\mathbf s$ là kích thước khung hiển thị. Quan hệ hiển thị là:
\begin{equation}
 \mathbf q_{\text{màn hình}}=z\mathbf q+\mathbf p,\qquad
 \mathbf p_{\text{mong muốn}}=\frac{\mathbf s}{2}-z\mathbf q.
\end{equation}
Độ pan sau đó được giới hạn để khung nhìn không lộ vùng trống ngoài nền bản đồ. Vì có giới hạn này, một tàu sát mép bản đồ có thể không nằm đúng tâm khung. Ví dụ khung rộng 400 đơn vị, marker ở $q_x=100$ và zoom $z=9$ cần $p_x=200-900=-700$ trước khi giới hạn biên.

\demopic[.86]{zoom_map.jpg}{Minimap phóng đại ở góc phải tập trung vào Player Ship và các Friendly AI, trong khi góc nhìn ngắm pháo vẫn được giữ cho gameplay.}{fig:minimapdetail}

M chỉ đổi cách xem bản đồ, không tự chuyển chuột khỏi ngắm bắn. Người chơi giữ Left Alt khi muốn dùng con trỏ để kéo hoặc cuộn bản đồ; thả Alt trả quyền nhìn/ngắm cho gameplay. Kéo thủ công tạm dừng việc tự theo để người dùng xem khu vực khác; tắt/bật M khôi phục theo tàu. Zoom camera pháo và zoom minimap là hai thao tác khác nhau.

Hai đội sử dụng texture marker riêng. Vì texture được nạp bằng đường dẫn chuỗi, thư mục \texttt{/Game/Game/UI/Minimap} được thêm vào danh sách cook để các icon có mặt trong bản đóng gói. Đường dẫn này tương ứng với \path{Content/Game/UI/Minimap}: \texttt{/Game} đầu là gốc asset đại diện cho Content, còn Game tiếp theo là thư mục con. Đây là một vị trí dữ liệu, không phải hai bản sao của minimap.


\section{Camera tàu, toàn cảnh và Full AI}
Camera theo tàu giúp quan sát chuyển động thân, thao tác ngắm và tương tác với mặt nước. Phím F mở góc nhìn toàn cảnh để đánh giá vị trí hai đội; nhấn lại để trở về góc gần. Phím Q đổi đối tượng đang theo dõi và đồng bộ HUD/minimap mà không làm mất AIController của tàu đó.

Trong Full AI, các slot đều do AI điều khiển và người dùng đóng vai trò quan sát. Camera tự do dùng W/S theo hướng nhìn, A/D dịch ngang, chuột đổi hướng và cuộn để thay tốc độ bay. Đây là chế độ phù hợp khi muốn quan sát một trận tự động hoặc kiểm tra nhiều tàu cùng đi tuyến.

Khi tàu người chơi chết, Q/Tab cho phép theo đồng đội còn sống; nếu không có đối tượng thích hợp thì trở về toàn cảnh. Hệ thống bảo vệ tham chiếu camera/HUD khi xác tàu được dọn, tránh giữ một actor đã bị hủy làm mục tiêu theo dõi.

\clearpage\section{Thời gian thực, theo lượt và chơi lại}
Real-Time cho các tàu đủ điều kiện hoạt động đồng thời. Turn-Based phân cửa sổ hành động 10 giây cho từng tàu và luân phiên giữa các đội, bỏ qua tàu đã chết. Chỉ tàu có lượt được phát lệnh chủ động theo quy tắc; mặt nước, vật lý và projectile đã tạo vẫn tiếp tục được mô phỏng. Vì vậy đây không phải cơ chế đóng băng toàn bộ thế giới giữa hai lượt.

Khi kết thúc trận, Match Summary cho biết bên thắng và thống kê từng tàu. \textbf{RETURN TO MAP SETTING} mở lại menu để thay đội hình hoặc bố trí. \textbf{RESTART} cho phép chơi lại qua luồng trận mới. \textbf{QUIT} kết thúc phiên chơi. Trở về xem trận đã kết thúc không hồi sinh các tàu đã bị tiêu diệt.

Một trận mới đặt lại roster, HP, đạn, bộ đếm, điểm số và các tham chiếu theo dõi thông qua vòng đời khởi tạo. Việc tách trở lại menu với tiếp tục quan sát giúp người dùng hiểu rõ mình đang xem kết quả cũ hay đã bắt đầu một trận khác.

\begin{figure}[H]\centering
\includegraphics[width=.95\linewidth,trim=460bp 450bp 450bp 450bp,clip]{images/result.jpg}
\caption{Khung Match Summary từ ảnh demo: DEFEAT cho đội người chơi, thống kê hai đội và các nút trở lại thiết lập, chơi lại, thoát.}\label{fig:result}
\end{figure}
\begingroup\setstretch{1.1}\setlength{\parskip}{3pt}
\chapter{Kết quả trình diễn và hướng phát triển}
\section{Tuần tra nhiều tàu và trận tự động}
Các bài chạy của dự án đã ghi nhận sáu tàu Medium đi tuần tra liên tục trong khoảng 1.350,93 giây mô phỏng, mỗi tàu có ít nhất bốn lượt tới đích. Với bài này, sát thương được tạm ngừng để quan sát đầy đủ vòng lặp chuyển động. Kết quả thể hiện khả năng phối hợp tuyến, đến đích và quay lại; không nên hiểu nó là thời lượng của một trận chiến tự nhiên.

Trong một trận 3v3 tự động, hai đội đã giao chiến và kết thúc với Blue thắng sau khoảng 803 giây mô phỏng. Cả sáu tàu đều phát bắn; ba tàu Red bị loại do HP cạn. Một tình huống trình diễn khác ở Hình~\ref{fig:result} cho thấy đội người chơi thất bại. Cả hai trường hợp sử dụng cùng luồng quản lý đội, thống kê và menu kết thúc trận.

Một kết quả chạy thành công có ý nghĩa với cấu hình tương ứng. Nó không bảo đảm mọi seed, loại tàu và bố trí đều phù hợp: tuyến quá hẹp hoặc điểm cuối sát bờ vẫn có thể không tạo được hành trình an toàn. Vì vậy, người trình diễn cần thiết kế và chạy thử chính kịch bản mình định sử dụng.

\section{Độ ổn định và chuyển động theo sóng}
Nhóm thử tàu Small ở các trạng thái đứng yên, hành trình, quay, lùi và vừa lùi vừa bắn. Góc roll/pitch được thu theo thời gian để phân biệt chuyển động sóng có kiểm soát với dao động quá mức.

Hai đại lượng được dùng để tóm tắt chuỗi góc $\theta_1,\ldots,\theta_n$ là:
\begin{equation}
 \theta_{RMS}=\sqrt{\frac1n\sum_{i=1}^{n}\theta_i^2},\qquad
 P_{95}=\operatorname{percentile}_{95}(|\theta_i|).
\end{equation}
RMS biểu diễn mức góc tổng thể theo bình phương trung bình. P95 cho biết 95\% mẫu có độ lớn không vượt giá trị đó; nó không phải góc lớn nhất. Hai chỉ số tính trên góc tuyệt đối của thân, vì vậy có thể bao gồm độ nghiêng cân bằng, không chỉ biên độ dao động quanh trung bình.

\begin{table}[H]\centering\small
\caption{Góc thân Small trong các trạng thái thử, đơn vị độ.}
\begin{tabular}{@{}lrrrr@{}}\toprule
\textbf{Trạng thái}&\textbf{Roll RMS}&\textbf{Pitch RMS}&\textbf{Roll P95}&\textbf{Pitch P95}\\\midrule
Đứng yên&0,069&1,117&0,125&2,034\\
Hành trình&0,053&0,706&0,075&0,970\\
Quay&0,354&0,569&0,597&0,969\\
Lùi&0,635&1,205&0,737&2,507\\
Lùi và bắn&0,653&1,179&0,793&2,374\\\bottomrule
\end{tabular}\end{table}

Kết quả cho thấy chuyển động không có cùng biên độ trên mọi trục: pitch rõ hơn roll ở các trạng thái trên. Một số đỉnh pitch nằm trong khoảng 2--3 độ, phù hợp mục tiêu giữ Small còn dập dềnh nhẹ. Khả năng rẽ cần được đánh giá cùng độ ổn định; một thân bị tăng damping tới mức gần đứng im về góc chưa chắc cho trải nghiệm điều khiển tốt.

\section{Những cải thiện về trải nghiệm trình diễn}
Chuyển động thân được hiệu chỉnh theo từng profile thay vì chỉ tăng damping chung. Lực cản theo trục giữ độ đầm khi ngắm nhưng để bánh lái tiếp tục tạo yaw. Tại bờ, lớp quét thân và cảnh báo nước nông giúp người chơi nhận biết hướng đi nguy hiểm và còn đường rẽ/lùi thoát ra.

Ở chiến đấu, thông báo tiêu diệt bổ sung tên đối tượng để người xem hiểu một tàu nào vừa bị loại. Bảng kết quả liên kết thống kê với các nút quay lại thiết lập, chơi lại và thoát. Các hình debug của projectile không còn được dùng làm hiệu ứng trúng đạn thông thường, giúp màn trình diễn tập trung vào hiệu ứng cháy/nổ và trạng thái sức bền.

Ở quan sát, minimap theo tàu giúp không phải tự tìm khu vực giao chiến khi phóng đại. Việc chỉ bật tương tác bản đồ lúc giữ Alt giữ được thao tác nhìn và bắn bằng chuột trong phần lớn thời gian chơi. Cảnh báo AMMO đỏ khi còn ít đạn giúp người chơi chủ động hơn trước thời gian nạp.

\section{Giới hạn và hướng phát triển}
Sản phẩm hiện phục vụ thử nghiệm và trình diễn trong phạm vi đội hình nhỏ. Những giới hạn chính gồm độ chính xác phụ thuộc dữ liệu terrain/collision, khoảng trống cần thiết khi tàu dài quay đầu, và hành vi nhiều phương tiện ở cửa hẹp. Bộ giải đạn đạo dùng xấp xỉ chuyển động mục tiêu trong khoảng dự báo ngắn; khi mục tiêu đổi hướng nhanh, sai lệch bắn đón vẫn có thể tăng.

Các thông số class là tham số hiệu chỉnh trong engine. Mô hình phao hình cầu chưa mô tả đầy đủ thủy động lực học, ngập khoang, sóng phức tạp hoặc đặc tính thật của từng loại tàu. Độ khó đã có cơ chế sai số và trễ nhưng cần thêm bài so sánh tỷ lệ trúng trong điều kiện tương đương. Các thống kê chạy trong Editor cũng chưa đủ để kết luận hiệu năng trên mọi máy.

Hướng phát triển tiếp theo là mở rộng kịch bản thử ở các cửa hẹp, cải thiện phối hợp tránh nhau, bổ sung chỉ số thời gian lập đường và tải CPU/GPU theo cảnh, đồng thời thử ổn định dài hạn với projectile và hiệu ứng. Có thể tiếp tục mở rộng mô hình trang bị, kịch bản nhiệm vụ và khả năng dùng nhiều bộ dữ liệu địa lý. Những hướng này cần được đánh giá trên từng mục tiêu cụ thể trước khi coi là tính năng đã hoàn thành.

\section{Kết luận}
Nhóm đã xây dựng môi trường hải chiến kết hợp bản đồ địa lý, phương tiện có quán tính, AI điều hướng và vòng đời chiến đấu. Người sử dụng có thể tạo đội hình, đặt tuyến A/B hoặc sinh bố trí từ seed, điều khiển một tàu hay quan sát toàn đội AI. Các chức năng vũ khí, sát thương, cháy/chìm, HUD và kết quả trận được kết nối trong cùng luồng sử dụng.

Giá trị kỹ thuật chính nằm ở sự phối hợp giữa hình học đường đi và chuyển động thực: A* cung cấp tuyến, bộ điều khiển biến tuyến thành ga/bánh lái, còn vật lý và cảm biến quyết định thân tàu thực sự đi như thế nào. Cách tổ chức này tạo cơ sở để tiếp tục phát triển các kịch bản mô phỏng trên địa hình thực tế.

\par\endgroup
\appendix
\chapter{Quy trình demo nhanh và xử lý tình huống}
\section{Một kịch bản từ thiết lập đến kết quả}
\begin{enumerate}
\item Mở level, đặt PlayerStart và các AI trong vùng nước. Gán A/B riêng, chừa vùng quay ở mỗi đầu tuyến rồi lưu.
\item Play và chọn bố trí thủ công. Đặt ba tàu mỗi bên, Player + AI, Medium cho người chơi, Normal và Real-Time.
\item Start New Match, chờ chuẩn bị hoàn tất. Thử tiến, hãm/lùi và rẽ hai phía; quan sát tốc độ và hướng trên HUD.
\item Giữ chuột phải để ngắm, bắn đối phương khi đủ điều kiện. Theo dõi HP, AMMO và đồng hồ nạp.
\item Nhấn M để xem gần đội hình; thử tiếp tục nhìn/ngắm. Giữ Alt để kéo bản đồ rồi thả để về gameplay.
\item Dùng Q và F quan sát các AI đang chạy A/B. Đọc kill feed khi một tàu bị tiêu diệt.
\item Khi có kết quả, mở lại thiết lập. Chuyển sang Full AI để quan sát trận tự động hoặc thay loại tàu và bắt đầu kịch bản khác.
\end{enumerate}

Nếu muốn trình diễn chiến thuật truy đuổi/rút lui, chuyển sang sinh bố trí tự động vì chế độ thủ công giữ tuyến A/B. Nếu muốn tập trung đo vòng lặp di chuyển, đặt hai đội xa nhau và theo dõi nhiều chặng trước khi giao chiến.

\clearpage\section{Các tình huống thường gặp}
\begin{longtable}{@{}p{35mm}p{110mm}@{}}
\caption{Hướng xử lý trong khi chuẩn bị và trình diễn.}\\
\toprule\textbf{Hiện tượng}&\textbf{Cách kiểm tra và xử lý}\\\midrule\endfirsthead
\toprule\textbf{Hiện tượng}&\textbf{Cách kiểm tra và xử lý}\\\midrule\endhead
Setup Failed&Kiểm tra đất/nước tại xuất phát và A/B, kết nối ba chặng, vùng quay và marker của slot Full AI. Đọc lý do lỗi rồi sửa bố trí trong Editor.\\
Một AI đứng yên&Kiểm tra đang pause, trận đã kết thúc, tàu đã chết, lượt hành động, trạng thái chờ hoặc tuyến bị chặn; xác nhận A/B của đúng actor.\\
Không quay lại B&A/B phải khác nhau và đủ xa; kiểm tra vùng đến, không gian quay và việc tàu đang giao chiến theo quyền chiến thuật hay giữ tuyến thủ công.\\
Không bắn được&Kiểm tra camera ngắm, tâm ngắm, băng đạn, reload/cooldown, lượt, góc pháo và vật cản. Đổi góc khỏi chính boong tàu rồi nhấn lại.\\
Bắn không mất HP&Thử đúng đối phương; kiểm tra hit có tới thân hay bị địa hình chắn. Self/friendly damage được chặn, xác tàu không còn là mục tiêu sống.\\
Tàu kênh gần đảo&Không đặt thân đè đất; giảm tốc, rẽ hoặc lùi khi có cảnh báo. Kiểm tra dữ liệu collision, không tắt toàn bộ collision của đảo.\\
Tàu có vẻ chậm&Phân biệt Cruise Speed với giới hạn tiến; AI có thể đang phanh, quay hoặc tránh. Thử trên đoạn nước mở trước khi tăng profile.\\
Chuột dùng cho map&Thả Left Alt để về gameplay. Trong Editor, Shift+F1 là cơ chế thả chuột của PIE và khác với Alt tương tác minimap.\\
Sửa map không được giữ&Dừng PIE, sửa ở Editor world rồi lưu. Actor sinh hoặc di chuyển trong phiên Play không tự trở thành bố trí xuất phát của lần sau.\\\bottomrule
\end{longtable}

\renewcommand{\bibname}{Tài liệu tham khảo}\cleardoublepage\phantomsection\addcontentsline{toc}{chapter}{Tài liệu tham khảo}
\begin{thebibliography}{9}
\bibitem{project} Nhóm phát triển. \textit{VDT Phase 2 Project: mã nguồn và tài nguyên mô phỏng hải chiến}. \url{https://github.com/Mr-Gwenchana/VDT-Phase2-Project}.
\bibitem{cesium} Cesium. \textit{ACesiumGeoreference Class Reference}. Tài liệu hệ tọa độ và chuyển đổi vị trí. \url{https://cesium.com/learn/cesium-unreal/ref-doc/classACesiumGeoreference.html}.
\bibitem{water} Epic Games. \textit{Water Buoyancy Component in Unreal Engine}. Nguyên lý xấp xỉ nổi bằng phao hình cầu. \url{https://dev.epicgames.com/documentation/en-us/unreal-engine/water-buoyancy-component-in-unreal-engine}.
\bibitem{input} Epic Games. \textit{Enhanced Input in Unreal Engine}. Cơ chế Input Action và Mapping Context. \url{https://dev.epicgames.com/documentation/en-us/unreal-engine/enhanced-input-in-unreal-engine}.
\bibitem{quantized} CesiumGS. \textit{Quantized-mesh terrain format}. Định dạng terrain và tệp mô tả layer.json. \url{https://github.com/CesiumGS/quantized-mesh/blob/main/SPECIFICATION.md}.
\bibitem{cesiumfaq} Cesium. \textit{Frequently Asked Questions: Data Formats}. Terrain quantized-mesh và lớp ảnh TMS từ dữ liệu cục bộ. \url{https://cesium.com/learn/unreal/unreal-faq/}.
\end{thebibliography}
\end{document}
